\documentclass[dvipdfmx,a4paper,oneside]{jsbook}
\usepackage[utf8]{inputenc}
\usepackage{amsmath}
\usepackage{amsfonts}
\usepackage{amssymb}
\usepackage{geometry}
\usepackage[dvipdfmx,pageanchor=false]{hyperref}
\usepackage{pxjahyper}
\usepackage{setspace}
\usepackage{ascmac}
\usepackage{booktabs}
\usepackage{enumitem}

\geometry{left=25mm,right=25mm,top=30mm,bottom=25mm}
\onehalfspacing

% 数式環境の調整
\numberwithin{equation}{chapter}

\title{クリフォード代数から射影幾何代数へ\\
― 二重時空理論のための \(G(3,1,1)\) 入門 ―}
\author{二重時空理論プロジェクト}
\date{}

\begin{document}

\maketitle
\tableofcontents
\newpage

%==================================================
\chapter*{序章　本書の目的と全体像}
\addcontentsline{toc}{chapter}{序章　本書の目的と全体像}

\section*{二重時空理論とは何か}

本書は、二重時空理論（Double Spacetime Theory, DST）を射影幾何代数 \(G(3,1,1)\) の言語で学ぶ入門書である。
DST の出発点は、時空が宇宙全体に一つだけ敷かれた連続体である、という従来の発想を改めることにある。
各粒子は、それぞれ\textbf{通常時空}と\textbf{双対時空}という二つのミンコフスキー構造を内在的に持つ。
通常時空は、私たちが日常で使う特殊相対性理論の舞台である。
双対時空は、コンパクト化された位相空間を持ち、時間の矢印が逆向きであるという特徴を持つ。

真空や自由落下の理想的な状態では、二つの時空は完全に同期している。
質量や運動があると、通常時空のローター \(R_{\mathrm{usual}}\) と双対時空のローター \(R_{\mathrm{dual}}\) の間にずれが生じる。
このずれこそが\textbf{ねじれ不整合}（torsional mismatch）であり、DST において重力と慣性はこの代数的な「角度のズレ」として記述される。
背景計量や曲率テンソルを最初から仮定するのではなく、粒子内在の双対構造から時空の効果を導く——これが DST の基本的な立場である。

本書で扱う PGA 版では、さらに第三の要素として\textbf{並進}が加わる。
並進は、空間や時間方向への実際の移動を司る自由度である。
従来の \(\mathrm{Cl}(3,1)\) だけでは並進は外付け的に与えられていたが、\(G(3,1,1)\) ではヌルベクトル \(e_4\) から生じる四つの生成子 \(N_\mu = e_4 \wedge e_\mu\) として代数の内部に内在化される。
したがって PGA 版 DST の構造は、\textbf{通常時空・双対時空・並進}の三者が協調する枠組みとなる。

三者の役割分担は次のように理解できる。
通常時空は、ローレンツ変換や光速での進行方向の「潰れ」といった、相対論の標準的現象を担う。
双対時空は、通常時空で失われがちな位相情報——とくに進行方向に関する情報——をコンパクトな位相として保存する。
並進セクターは、保存された情報をもとに、実際の位置のずれを光的な移動として実行する。
この三者構造こそが、本書全体を貫く物理的・代数的な骨格である。

\section*{なぜ \(G(3,1,1)\) なのか}

DST の親理論は、ミンコフスキー時空の幾何代数 \(\mathrm{Cl}(3,1)\) 上に構築される。
この代数は双四元数と同型であり、16 次元の構造を持つ。
四つのベクトル生成子 \(e_0 = j,\ e_1 = kI,\ e_2 = kJ,\ e_3 = kK\) は符号 \((-,+,+,+)\) のミンコフスキー内積を満たし、ローレンツ群の作用をローターとして表現する。
ここまでで、回転とブースト——すなわち \(\mathfrak{so}(3,1)\) の六自由度——は自然に扱える。

しかし \(\mathrm{Cl}(3,1)\) だけでは、時空並進は代数の内部からは生まれない。
並進を幾何的に統合するため、本書は射影幾何代数 \(G(3,1,1)\) へ理論を持ち上げる。
具体的には、平方がゼロの\textbf{ヌルベクトル} \(e_4\)（\(e_4^2 = 0\)）を一つ付加する。
\(e_4\) は既存の四ベクトルと反可換し（\(\{e_4, e_\mu\} = 0\)）、双四元数の擬スカラー \(i = e_0 e_1 e_2 e_3\) とは可換である。
この性質により、親理論の双対写像 \(X \mapsto Xi\) は \(\mathrm{Cl}(3,1)\) 部分代数上で従来どおり作用しつつ、新しいヌル方向 \(e_4\) は独立に保たれる。

\(G(3,1,1)\) のバイベクトル空間は \(\binom{5}{2} = 10\) 次元である。
拡張されたねじれ不整合の生成子は、この 10 次元空間に次の三つに分類される。

\begin{itemize}[leftmargin=2em]
  \item \textbf{双曲生成子}（3 個、平方 \(+1\)）：\(iI = e_0 e_1,\ iJ = e_0 e_2,\ iK = e_0 e_3\)。
  通常時空側のブーストを担い、\(\mathfrak{so}(3,1)\) の双曲的部分に対応する。
  \item \textbf{巡回生成子}（3 個、平方 \(-1\)）：\(I = e_3 e_2,\ J = e_1 e_3,\ K = e_2 e_1\)。
  双対時空側の回転を担い、\(\mathfrak{so}(3,1)\) のもう一つのコピーに対応する。
  \item \textbf{ヌル生成子}（4 個、平方 \(0\)）：\(N_\mu = e_4 \wedge e_\mu\)（\(\mu = 0,1,2,3\)）。
  並進を担い、可換な理想 \(\{P_\mu\}\) を形成する。
\end{itemize}

この 10 次元構造は、ポアンカレ代数 \(\mathfrak{iso}(3,1)\) のリー代数と一致する。
六つの双曲・巡回生成子が \(\mathfrak{so}(3,1) \oplus \mathfrak{so}(3,1)\) を張り、四つのヌル生成子が並進理想を与える。
並進がゲージ場や座標変換として外付けされるのではなく、代数のヌルセクターから有限のシフトとして現れる——これが \(G(3,1,1)\) を選ぶ中心的な理由である。

ヌル生成子の代数的性質 \(N_\mu N_\nu = 0\)（すべての \(\mu, \nu\) に対して）は、並進の指数写像を一次多項式に打ち切る。
すなわち
\[
\exp\!\Bigl(\tfrac{1}{2}\lambda^\mu N_\mu\Bigr)
= 1 + \tfrac{1}{2}\lambda^\mu N_\mu
\]
となり、並進は高次補正なしに有限の移動として実現される。
拡張モーター \(M = \exp(\Omega_{\mathrm{biv}}^{(5)})\) はさらに反転奇性 \(\widetilde{N_\mu} = -N_\mu\) によりユニタリ（\(M\widetilde{M} = 1\)）となり、因果的な伝播を保つ。
これらの詳細は第 III 部で体系的に扱う。

\section*{ニュートン‐ウィグナー問題を解くための数学的道具立て}

相対論的量子力学において、質量ゼロ粒子——とくに光子——の位置を定義することは長年の難問であった。
1949 年の Newton と Wigner の論文以降、\textbf{ニュートン‐ウィグナー問題}として知られるこの問題は、質量のある粒子では共変な位置演算子が構成できる一方、光子では進行方向の自由度がローレンツ収縮によって「潰れ」、一意な位置を定義できないことを示している。

問題の核心は次の二点に集約される。
第一に、光子が光速で進むとき、進行方向の位置座標は観測者のローレンツ変換に依存して消えてしまう。
第二に、失われた進行方向の情報を別の場所に保存する仕組みが、従来の単一時空の枠組みにはなかった。

PGA 版 DST は、この問題を三者構造で解く。
通常時空で潰れた情報は双対時空のコンパクト位相に移され、実際の移動は並進セクターが光的に実行する。
光子では通常時空と双対時空のローターが\textbf{完全反同期}となり、ねじれ不整合がゼロに固定される。
この条件のもとで並進パラメータ \(\lambda^\mu\) は光円錐上に強制され（\(\eta_{\mu\nu}\lambda^\mu\lambda^\nu = 0\)）、光子の移動は純粋な光速並進となる。

本書で学ぶ数学的道具と、物理問題の対応を表 \ref{tab:tools-nw} に示す。
ここでは対応関係の概観のみを示し、厳密な構成と証明は第 IV 部・第 V 部に委ねる。

\begin{table}[htbp]
  \centering
  \caption{ニュートン‐ウィグナー問題と本書の道具立て}
  \label{tab:tools-nw}
  \begin{tabular}{@{}lll@{}}
    \toprule
    物理的要請 & 本書の道具 & 登場部 \\
    \midrule
    回転・ブーストの統一 & ローター・ベルソル & 第 I--II 部 \\
    並進の内在化 & \(e_4\), \(N_\mu\), モーター & 第 III 部 \\
    進行方向情報の保存 & 双対時空のコンパクト位相 & 第 IV 部 \\
    光的移動の強制 & ねじれ不整合ゼロ \(\Rightarrow\) 光円錐条件 & 第 IV--V 部 \\
    位置の一意性 & 横方向 \(+\) 位相 \(+\) 光的並進 & 第 V 部 \\
    \bottomrule
  \end{tabular}
\end{table}

位置演算子は、通常時空の進行方向に垂直な横方向座標、双対時空の進行方向位相、および並進セクターによる光的移動の三つを組み合わせて構成される。
この構成により、すべての観測者で一致する共変な位置演算子が得られ、70 年以上にわたる問題が解決される。
本書は、この解決のための代数を段階的に積み上げることを目的とする。

\section*{本書の前提知識と記法規約}

\subsection*{前提知識}

本書は一般大学生を想定する。
以下の内容は既知とみなし、詳しい復習は行わない。

\begin{itemize}[leftmargin=2em]
  \item 線形代数（ベクトル空間、基底、内積の基本）
  \item 三次元ユークリッド幾何（内積、外積、回転の直感）
  \item 四元数の初歩（虚数単位 \(I, J, K\) と回転の関係）
  \item 特殊相対性理論の初歩（ローレンツ変換、光速度不変、\(E = mc^2\) の直感）
\end{itemize}

クリフォード代数、射影幾何代数、DST の物理的内容は、本書の第 I 部から順に導入する。
必要に応じて付録も参照されたい。

\subsection*{用語と記法}

本書では、幾何代数・DST 文献で用いられる英語用語と日本語訳を表 \ref{tab:terminology} のとおり対応させる。

\begin{table}[htbp]
  \centering
  \caption{用語対応表}
  \label{tab:terminology}
  \begin{tabular}{@{}ll@{}}
    \toprule
    英語 & 日本語 \\
    \midrule
    dual spacetime & 双対時空 \\
    biquaternion & 双四元数 \\
    vector & ベクトル \\
    bivector & バイベクトル \\
    blade & ブレード \\
    multivector & 多重ベクトル \\
    grade & 階数 \\
    reversion & 反転 \\
    versor & ベルソル \\
    rotor & ローター \\
    torsion & トーション \\
    torsion star & トーション星 \\
    torsional mismatch & ねじれ不整合 \\
    Killing form & キリング形式 \\
    \bottomrule
  \end{tabular}
\end{table}

\subsection*{符号規約}

通常時空の四ベクトルはミンコフスキー符号 \((-,+,+,+)\) を用いる。
すなわち \(e_0^2 = -1\)、\(e_1^2 = e_2^2 = e_3^2 = +1\)、異なる \(e_\mu\) は反可換である。
付加するヌルベクトルは \(e_4^2 = 0\) かつ \(\{e_4, e_\mu\} = 0\) を満たす。
幾何積は、ベクトルの順序を保持する（非可換）積とする。
反転（reversion）は、バイベクトル以上の反転奇性を持つ元に対して符号を反転させる操作である。
詳細な符号一覧は付録 A にまとめる。

\subsection*{本書の構成}

本書は五部構成である。
第 I 部では四元数から出発し、クリフォード代数の公理、多重ベクトル、ベルソルとローターを学ぶ。
第 II 部では時空代数 \(G(3,1)\) を導入し、ローレンツ変換と電磁気の幾何代数的記述に触れる。
第 III 部ではヌルベクトル \(e_4\) を加えた \(G(3,1,1)\) を構築し、並進・モーター・双対性の代数的基盤を整える。
第 IV 部では通常時空と双対時空の二重構造、ねじれ不整合、位置演算子の構成へと DST の数学的枠組みを完成させる。
第 V 部ではニュートン‐ウィグナー問題の解決と、重力・慣性への拡張の展望を示す。

各章末には演習問題を置く。
手計算による幾何積の練習を通じて、代数の操作感覚を身につけることが本書の学習目標である。

\newpage

%==================================================
\part{第I部　クリフォード代数の基礎}

\chapter{四元数からクリフォード代数へ}

序章で述べた \(G(3,1,1)\) や双四元数は、すべてクリフォード代数という枠組みの上に立つ。
本章では、その出発点となる四元数を復習し、幾何積という一つの積に内積と外積を統合する考え方へと持ち上げる。
第 II 部以降で現れるバイベクトル生成子 \(\{I,J,K\}\) やローターは、ここで整えた言語の自然な延長である。

\section{四元数の復習（回転としての解釈）}

四元数は、3 次元空間の回転を代数の積だけで記述する道具である。
本書では論文の巡回生成子と記法を揃えるため、虚数単位を大文字 \(I,J,K\) と書く（従来文献では \(i,j,k\) と書かれることも多い）。
四元数 \(q\) は
\begin{equation}
  q = w + xI + yJ + zK, \qquad w,x,y,z \in \mathbb{R}
\end{equation}
と表される。
乗法規則は
\begin{align}
  IJ &= K, & JI &= -K, \notag \\
  JK &= I, & KJ &= -I, \notag \\
  KI &= J, & IK &= -J, \notag \\
  I^2 &= J^2 = K^2 = IJK = -1
  \label{eq:quat-mult}
\end{align}
である。
非可換性 \(IJ \neq JI\) が、回転の合成が順序に依存することと対応する。

共役は \(\widetilde{q} = w - xI - yJ - zK\)、ノルムは \(|q|^2 = q\widetilde{q} = w^2 + x^2 + y^2 + z^2\) と定義する。
\(|q|=1\) となる単位四元数が回転を表現する。

純虚四元数 \(u = xI + yJ + zK\) の平方は \(u^2 = -|u|^2\) である。
単位純虚四元数 \(\hat{q}\)（\(|\hat{q}|=1\)）は回転軸を表し、オイラーの公式の拡張
\begin{equation}
  e^{\hat{q}\,\theta/2} = \cos\frac{\theta}{2} + \hat{q}\,\sin\frac{\theta}{2}
  \label{eq:quat-rotor}
\end{equation}
は単位四元数——すなわち\textbf{ローター}——を与える。
指数の引数に半角 \(\theta/2\) が現れるのは、回転角 \(\theta\) と区別するためである。

3 次元ベクトル \(\mathbf{v}=(x,y,z)\) を純虚四元数 \(v = xI + yJ + zK\) と同一視すると、軸 \(\hat{q}\) 周りの角度 \(\theta\) の回転は
\begin{equation}
  v' = q\,v\,\widetilde{q}
  \label{eq:quat-sandwich}
\end{equation}
の\textbf{サンドイッチ積}で与えられる。
回転の合成は四元数の積 \(q_2 q_1\) として表され、ジンバルロックのような特異点を避けられる。

四元数が示す核心は、「幾何変換が代数の元として書ける」という点にある。
後の章で現れる巡回生成子 \(I=e_3e_2,\ J=e_1e_3,\ K=e_2e_1\) も、平方 \(-1\) を持つ同型の対象として、この代数の延長線上に位置する。

\section{幾何積の公理的導入}

四元数の乗法は、実部と虚部の内積・外積を一つの積にまとめたものとして読める。
この統合を一般化したのが\textbf{クリフォード代数}である。

\(n\) 次元ベクトル空間 \(V\) 上に、結合的で双線形な積——\textbf{幾何積}——を導入する。
任意のベクトル \(a \in V\) に対し、二次形式 \(Q(a)\) を用いて公理
\begin{equation}
  a^2 = Q(a)
  \label{eq:clifford-axiom}
\end{equation}
を課す。
直交基底 \(\{e_1,\ldots,e_n\}\) を取れば、クリフォード関係
\begin{equation}
  e_i e_j + e_j e_i = 2\,\eta_{ij}
  \label{eq:clifford-relation}
\end{equation}
が得られる。
ここで \(\eta_{ij}\) は基底間の内積（本章ではまずユークリッド空間 \(\eta_{ij}=\delta_{ij}\) を想定する）。

幾何積 \(ab\) は一般に可換でも反可換でもない。
しかし任意の \(a,b\) を対称部分と反対称部分に分解すれば、ベクトル解析で知られる\textbf{内積}と\textbf{外積}が現れる（次節）。
これが「一つの積から幾何が生まれる」というクリフォード代数の出発点である。

\(n\) 個の基底ベクトルから出発すると、スカラー（0 次）、ベクトル（1 次）、バイベクトル（2 次）、\(\ldots\)、擬スカラー（\(n\) 次）まで、合計 \(2^n\) 次元の代数が生成される。
この全体を\textbf{幾何代数} \(G(n)\) と呼ぶ。
\(n\) 次元空間のユークリッド幾何代数を \(G(n)\) と書き、符号付きの場合は \(G(p,q)\) と書く（第 II 部で \(G(3,1)\) が現れる）。
多重ベクトルの階数分解や操作の詳細は次章に委ねる。

\section{内積・外積の分解とクリフォード関係}

ベクトル \(a,b\) に対し、幾何積は次のように分解される：
\begin{equation}
  ab = a\cdot b + a\wedge b,
  \label{eq:gp-decomp}
\end{equation}
\begin{equation}
  a\cdot b = \tfrac{1}{2}(ab + ba), \qquad
  a\wedge b = \tfrac{1}{2}(ab - ba).
  \label{eq:dot-wedge}
\end{equation}
\(a\cdot b\) はスカラー（0 次）、\(a\wedge b\) は\textbf{バイベクトル}（2 次）である。

この分解から、計算上の基本ルールが読み取れる。
\begin{itemize}[leftmargin=2em]
  \item \(a \parallel b\)（平行）なら \(a\wedge b = 0\)。幾何積は内積のみになる。
  \item \(a \perp b\)（直交）なら \(a\cdot b = 0\)。幾何積は外積のみになる：\(ab = a\wedge b\)。
  \item バイベクトル \(a\wedge b\) は、\(a\) と \(b\) が張る\textbf{向き付き平面}を表す。
  面積の大きさは \(|a\wedge b| = |a|\,|b|\,|\sin\theta|\)（\(\theta\) はなす角）である。
\end{itemize}
後の章で現れる 10 次元バイベクトル空間も、この「平面要素」の高次元版として理解できる。

具体例を二つ示す。
直交単位ベクトル \(e_1,e_2\) に対し
\begin{equation}
  e_1 e_2 = e_1 \wedge e_2, \qquad (e_1 e_2)^2 = e_1 e_2 e_1 e_2 = -e_1^2 e_2^2 = -1.
  \label{eq:bivector-square}
\end{equation}
バイベクトルの平方が \(-1\) になるのは、平面内の「90 度回転」を内包しているためである。
また \(a = e_1 + e_2\) とすると
\begin{equation}
  a^2 = (e_1 + e_2)^2 = e_1^2 + e_1 e_2 + e_2 e_1 + e_2^2
       = 1 + 0 + 1 = 2,
\end{equation}
となり、公理 \eqref{eq:clifford-axiom} の \(Q(a) = |a|^2 = 2\) と一致する。

\section{\texorpdfstring{低次元の例：\(G(2)\) と \(G(3)\)}{低次元の例：G(2) と G(3)}}

幾何代数の具体像を得るため、次元 2 と 3 の例を見る。

\paragraph{\(G(2)\)：平面の幾何代数}
\(e_1^2 = e_2^2 = 1\)、\(e_1 e_2 = -e_2 e_1\) なる直交基底 \(\{e_1, e_2\}\) から、4 次元代数
\[
G(2) = \operatorname{span}\{1,\; e_1,\; e_2,\; e_1 e_2\}
\]
が生成される。
擬スカラー \(I_2 := e_1 e_2\) は \(I_2^2 = -1\) を満たし、\(G(2)\) は複素数 \(\mathbb{C}\) と同型である。
平面内の回転は
\begin{equation}
  R(\theta) = e^{I_2 \theta} = \cos\theta + I_2 \sin\theta
\end{equation}
と書ける。
これは \eqref{eq:quat-rotor} の 2 次元版である。

\paragraph{\(G(3)\)：3 次元空間の幾何代数}
\(e_1^2 = e_2^2 = e_3^2 = 1\) なる直交基底から 8 次元代数 \(G(3)\) が生成される。
その\textbf{偶数部分代数}（スカラーとバイベクトルのみからなる部分代数）は 4 次元であり、四元数と同型である。
具体的な同一視は
\begin{equation}
  I = e_3 e_2, \qquad J = e_1 e_3, \qquad K = e_2 e_1
  \label{eq:quat-embedding}
\end{equation}
である。
計算すれば \(I^2 = J^2 = K^2 = -1\)、\(IJK = -1\) が成り立ち、第 1 節の四元数関係と一致する。

したがって四元数は、\(G(3)\) の偶数部分に他ならない。
3 次元回転のローター \eqref{eq:quat-rotor} も、バイベクトル \(B = \tfrac{1}{2}(xI + yJ + zK)\) を用いた \(e^B\) として幾何代数の言葉で書ける。
サンドイッチ積 \eqref{eq:quat-sandwich} も、次章以降で一般化される\textbf{ベルソル}作用の特殊例である。

\paragraph{第 II 部への接続}
符号を \((-,+,+,+)\) に変えた \(G(3,1)\) は 16 次元となり、双四元数 \(\mathrm{Cl}(3,1)\) と同型である（序章参照）。
さらにヌルベクトル \(e_4\) を加えた \(G(3,1,1)\) が第 III 部の舞台となる。
本章で学んだ「幾何積 → バイベクトル → ローター」の流れは、そのまま時空代数と PGA へと拡張される。

\section*{演習問題}

\begin{enumerate}[leftmargin=2em]
  \item \(I,J,K\) の乗法表 \eqref{eq:quat-mult} のうち、\(JK = I\) と \(KJ = -I\) を直接計算で確かめよ（\eqref{eq:quat-embedding} の定義を使ってもよい）。

  \item \(a = 3e_1\)、\(b = 4e_2\) のとき、\(ab\) を \eqref{eq:gp-decomp} に従って \(a\cdot b + a\wedge b\) に分解せよ。

  \item \(G(2)\) において \(e_1 e_2 e_1\) を計算せよ。
  結果が \(-e_2\) となることを確認し、\(e_1\) による反射の萌芽として解釈せよ。

  \item \eqref{eq:quat-embedding} の \(I,J,K\) が \(I^2 = J^2 = K^2 = -1\) および \(IJ = K\) を満たすことを確かめよ。

  \item （発展）直交する単位ベクトル \(e_i, e_j\)（\(i \neq j\)）に対し、\((e_i e_j)^2 = -1\) となる理由を、\eqref{eq:clifford-relation} と \eqref{eq:bivector-square} を用いて説明せよ。
\end{enumerate}

\chapter{多重ベクトルと階数}

前章では、幾何積の公理から \(G(2)\) と \(G(3)\) の具体例までを見た。
本章では、代数の元を\textbf{階数}（grade）で分類し、\textbf{多重ベクトル}として体系的に扱う。
第 II 部以降の \(G(3,1)\)、第 III 部の \(G(3,1,1)\) も同じ枠組みの高次元版である。
特に後者では 10 次元のバイベクトル空間や 3-ブレードによる平面表現が中心的役割を担うが、その代数的基礎はここで整える。

\section{ブレードと多重ベクトル}

幾何代数の元は、一次独立なベクトルの外積から生まれる\textbf{ブレード}（blade）を基本単位として理解できる。
\(k\) 個の一次独立なベクトル \(v_1,\ldots,v_k\) に対し、
\begin{equation}
  v_1 \wedge v_2 \wedge \cdots \wedge v_k
  \label{eq:k-blade}
\end{equation}
を \textbf{\(k\)-ブレード}と呼ぶ。
これは向き付き \(k\) 平行体——点（\(k=0\)）、有向線分（\(k=1\)）、有向平面（\(k=2\)）、有向体積（\(k=3\)）——を表す。
外積の順序を入れ替えると符号が反転する（\(a\wedge b = -b\wedge a\)）ため、ブレードは向きを持つ。

\textbf{多重ベクトル}（multivector）\(M\) は、各階数の成分を足し合わせた一般の元である：
\begin{equation}
  M = \sum_{k=0}^{n} \langle M \rangle_k,
  \label{eq:mv-decomp}
\end{equation}
ここで \(\langle M \rangle_k\) は \(M\) の\textbf{\(k\) 次部分}（grade-\(k\) part）を表す。
\(n\) 次元ベクトル空間から生成される幾何代数 \(G(n)\) は \(2^n\) 次元であり、階数 \(k\) の部分空間の次元は \(\binom{n}{k}\) である。

\paragraph{\(G(2)\) の基底と階数}
第 1 章で見た \(G(2)\) を階数の言葉で整理する。
直交基底 \(\{e_1,e_2\}\) から、4 つの線形独立な元が得られる：
\begin{center}
\begin{tabular}{@{}cll@{}}
  \toprule
  階数 & 基底元 & 幾何的意味 \\
  \midrule
  0 & \(1\) & スカラー \\
  1 & \(e_1,\ e_2\) & ベクトル \\
  2 & \(e_1 e_2\) & バイベクトル（有向面積） \\
  \bottomrule
\end{tabular}
\end{center}
擬スカラー \(I_2 := e_1 e_2\) は 2 次の最高階成分であり、\(I_2^2 = -1\) を満たす。

\paragraph{\(G(3)\) の基底と階数}
\(G(3)\) は 8 次元である：
\begin{center}
\begin{tabular}{@{}cll@{}}
  \toprule
  階数 & 基底元 & 幾何的意味 \\
  \midrule
  0 & \(1\) & スカラー \\
  1 & \(e_1,\ e_2,\ e_3\) & ベクトル \\
  2 & \(e_1 e_2,\ e_2 e_3,\ e_3 e_1\) & バイベクトル（有向平面） \\
  3 & \(e_1 e_2 e_3\) & 三ベクトル（有向体積）＝擬スカラー \\
  \bottomrule
\end{tabular}
\end{center}
第 1 章の四元数同一視 \eqref{eq:quat-embedding} では、偶数階数部分（スカラー＋バイベクトル、計 4 次元）だけが四元数代数をなす。
一方、ベクトル（1 次）と三ベクトル（3 次）は奇数階数部分を形成し、回転のローター単独では生成されない——この偶奇の区別は次章のベルソル論で重要になる。

\paragraph{第 III 部への予告}
\(G(3,1,1)\) では 5 次元ベクトル空間から 32 次元の代数が生成される。
とくにバイベクトル空間は \(\binom{5}{2}=10\) 次元であり、双曲・巡回・ヌルの三種類の生成子を収容する（序章参照）。
また射影幾何では\textbf{平面}が 3-ブレードとして表される。
例えば \(e_1\wedge e_2\wedge e_3\) は 3 次元部分空間内の有向体積要素であり、\(e_4\wedge e_1\wedge e_2\) のように \(e_4\) を含む 3-ブレードはヌル方向を持つ平面を表す。
これらの操作の詳細は第 III 部で扱う。

\section{階数分解}

\subsection{階数射影の定義}

任意の多重ベクトル \(M\) に対し、\textbf{階数射影} \(\langle\cdot\rangle_k\) は \(M\) から \(k\) 次部分だけを取り出す線形写像である。
次の性質を満たす：
\begin{align}
  \langle \alpha M + \beta N \rangle_k &= \alpha\langle M \rangle_k + \beta\langle N \rangle_k,
  \label{eq:grade-linear} \\
  \langle \langle M \rangle_k \rangle_\ell &= \delta_{k\ell}\,\langle M \rangle_k.
  \label{eq:grade-idempotent}
\end{align}
すなわち、一度 \(k\) 次部分を取り出せば、再度射影しても変わらない（冪等性）。
直交基底 \(\{e_1,\ldots,e_n\}\) では、\(k\) 個の添字を持つ基底積が \(\langle\cdot\rangle_k\) の像を張る。

例として \(M = 3 + 2e_1 + e_1 e_2 + 4e_1 e_2 e_3 \in G(3)\) を考える。
\begin{align}
  \langle M \rangle_0 &= 3, &
  \langle M \rangle_1 &= 2e_1, &
  \langle M \rangle_2 &= e_1 e_2, &
  \langle M \rangle_3 &= 4e_1 e_2 e_3.
\end{align}
\(\langle M \rangle_0 + \langle M \rangle_1 + \langle M \rangle_2 + \langle M \rangle_3 = M\) が成り立つ。

\subsection{幾何積の階数内容}

同じ階数 \(r\) の多重ベクトル \(A_r\) と \(s\) 次の \(B_s\) の幾何積 \(A_r B_s\) は、一般に複数の階数を含む。
可能な階数は
\begin{equation}
  \{|r-s|,\ |r-s|+2,\ \ldots,\ r+s\}
  \label{eq:grade-range}
\end{equation}
に限られる。
直交基底での計算から、この範囲外の成分は現れない。

この事実は、内積と外積の一般化を与える。
\(r\) 次の \(A_r\) と \(s\) 次の \(B_s\) に対し、
\begin{equation}
  A_r \cdot B_s = \langle A_r B_s \rangle_{|r-s|}, \qquad
  A_r \wedge B_s = \langle A_r B_s \rangle_{r+s}
  \label{eq:dot-wedge-general}
\end{equation}
と定義する。
\(r=s=1\) のとき \eqref{eq:dot-wedge-general} は前章の \eqref{eq:dot-wedge} に一致する。

\paragraph{計算例}
直交単位ベクトル \(e_1,e_2,e_3\) に対し、
\begin{equation}
  (e_1 e_2)^2 = e_1 e_2 e_1 e_2 = -e_1^2 e_2^2 = -1,
  \label{eq:biv-square-g3}
\end{equation}
これは 0 次部分のみである。
また三ベクトル \(I_3 := e_1 e_2 e_3\) の平方は
\begin{equation}
  I_3^2 = e_1 e_2 e_3 e_1 e_2 e_3 = (-1)^3 e_1^2 e_2^2 e_3^2 = -1.
  \label{eq:trivector-square}
\end{equation}

2 次同士の積の例として、\(A_2 = e_1 e_2\)、\(B_2 = e_2 e_3\) を計算する。
\begin{align}
  A_2 B_2 &= e_1 e_2 e_2 e_3 = e_1 e_2^2 e_3 = e_1 e_3, \\
  \langle A_2 B_2 \rangle_0 &= 0, \qquad
  \langle A_2 B_2 \rangle_2 = e_1 e_3, \qquad
  \langle A_2 B_2 \rangle_4 = 0.
\end{align}
\eqref{eq:grade-range} では \(r=s=2\) なので階数 \(\{0,2,4\}\) が可能だが、\(G(3)\) には 4 次成分は存在せず、実際には 0 次と 2 次のみが現れる。

\paragraph{高次元での階数分解}
5 次元の \(G(3,1,1)\) では、3 次同士の積 \(A_3 B_3\) の階数内容が
\begin{equation}
  A_3 B_3 = \langle A_3 B_3 \rangle_0 + \langle A_3 B_3 \rangle_2 + \langle A_3 B_3 \rangle_4
  \label{eq:grade-3-product}
\end{equation}
に限定される（6 次成分は存在しない）。
交換子 \(\tfrac{1}{2}(AB-BA)\) は 2 次部分を、反交換子 \(\tfrac{1}{2}(AB+BA)\) は 0 次と 4 次部分を抽出する——この性質は第 III--IV 部の平面幾何と双対時空の対応で決定的な役割を果たす。

\section{逆元・共役・ノルム}

\subsection{反転}

\textbf{反転}（reversion）は、幾何積におけるベクトルの順序を逆にする操作である。
\(k\)-ブレード \(v_1\wedge\cdots\wedge v_k\) に対し
\begin{equation}
  \widetilde{v_1 \wedge \cdots \wedge v_k} = v_k \wedge \cdots \wedge v_1
  \label{eq:reversion-blade}
\end{equation}
と定義し、線形性により一般の多重ベクトルへ拡張する。
\(k\)-ブレードに対する符号は
\begin{equation}
  \widetilde{A_k} = (-1)^{k(k-1)/2}\,A_k
  \label{eq:reversion-sign}
\end{equation}
である。
具体的には、スカラーはそのまま、ベクトルはそのまま、バイベクトルは符号反転、三ベクトルはそのまま、というパターンになる。

直交基底では \((e_i e_j)^{\sim} = e_j e_i = -e_i e_j\)（\(i\neq j\)）であるから、すべてのバイベクトルは\textbf{反転奇性}（reversal-odd）を持つ：
\begin{equation}
  \widetilde{e_i e_j} = -e_i e_j.
  \label{eq:biv-rev-odd}
\end{equation}
第 III 部で現れるヌル生成子 \(N_\mu = e_4\wedge e_\mu\) もバイベクトルであるため \(\widetilde{N_\mu} = -N_\mu\) となる。
これは並進を\textbf{ヌルベクトル}（1 次）ではなく\textbf{ヌルバイベクトル}（2 次）として実装する代数的理由の一つである——1 次のヌルベクトルは反転偶性を持ち、後述のユニタリ性 \(M\widetilde{M}=1\) が成り立たなくなる。

\subsection{階数対合}

\textbf{階数対合}（grade involution）\(\widehat{M}\) は、\(k\) 次部分に符号 \((-1)^k\) を付ける操作：
\begin{equation}
  \widehat{\langle M \rangle_k} = (-1)^k \langle M \rangle_k.
  \label{eq:grade-involution}
\end{equation}
偶数階数部分（スカラー、バイベクトル、\(\ldots\)）と奇数階数部分（ベクトル、三ベクトル、\(\ldots\)）を入れ替える符号を担う。
\(G(3)\) の偶数部分代数（四元数）や、次章で扱うベルソルの偶奇性の議論に現れる。

\subsection{ノルムと可逆性}

ユークリッド設定では、多重ベクトル \(M\) の\textbf{スカラー部ノルム}を
\begin{equation}
  \|M\|^2 = \langle M\widetilde{M}\rangle_0
  \label{eq:mv-norm}
\end{equation}
と定義する。
\(\|M\|^2\) がゼロでない実数のとき \(M\) は可逆であり、
\begin{equation}
  M^{-1} = \frac{\widetilde{M}}{\|M\|^2}
  \label{eq:mv-inverse}
\end{equation}
が成り立つ。
符号付き代数や退化代数（\(e_4^2=0\) など）では \eqref{eq:mv-norm} の解釈に注意が必要だが、基本的な操作原理は同じである。

\paragraph{例：単位四元数とローター}
\(G(3)\) の偶数部分 \(q = w + xI + yJ + zK\)（\(I=e_3e_2\) 等）に対し、
\begin{equation}
  q\widetilde{q} = w^2 + x^2 + y^2 + z^2 = |q|^2
\end{equation}
はスカラーである。
\(|q|=1\) なら \(q^{-1} = \widetilde{q}\) となり、第 1 章のローター \eqref{eq:quat-rotor} は
\begin{equation}
  R\widetilde{R} = 1
  \label{eq:rotor-unitarity}
\end{equation}
を満たす——これが回転を保つ代数的根拠である。

\section{擬スカラーと双対}

\subsection{擬スカラー}

\(n\) 次元ベクトル空間の直交基底 \(\{e_1,\ldots,e_n\}\) から、最高階数の基底元
\begin{equation}
  I_n := e_1 e_2 \cdots e_n
  \label{eq:pseudoscalar-def}
\end{equation}
を \textbf{擬スカラー}（pseudoscalar）と呼ぶ。
\(I_n\) は \(n\) 次ブレードであり、向き付き \(n\) 次元体積要素を表す。
符号付き代数では
\begin{equation}
  I_n^2 = (-1)^{n(n-1)/2}\,\eta_{1}\eta_{2}\cdots\eta_{n},
  \label{eq:pseudo-square}
\end{equation}
ここで \(\eta_i = e_i^2\) は各基底ベクトルの符号である。
ユークリッド \(G(3)\) では \(I_3^2 = -1\)（\eqref{eq:trivector-square}）、\(G(2)\) では \(I_2^2 = -1\) である。

\subsection{代数的双対}

擬スカラー \(I_n\) を用い、\(k\) 次の \(A_k\) の\textbf{代数的双対}（algebraic dual）を
\begin{equation}
  A_k^* := A_k I_n^{-1}
  \label{eq:algebraic-dual}
\end{equation}
と定義する（\(I_n^2 = \pm 1\) なら \(I_n^{-1} = \pm I_n\)）。
この操作は \(k\) 次と \((n-k)\) 次を対応させ、外積代数におけるホッジ双対 \(\star\) と一致する。

\paragraph{\(G(3)\) での例}
\(I_3 = e_1 e_2 e_3\)、\(I_3^{-1} = -I_3\) として、ベクトル \(e_1\) の双対は
\begin{equation}
  e_1^* = e_1 I_3^{-1} = e_1(-e_1 e_2 e_3) = -e_2 e_3,
\end{equation}
すなわち \(e_1\) に垂直な平面 \(e_2\wedge e_3\) に対応する（向きの符号は \(I_3^{-1}\) の取り方に依存する）。
逆に \((e_2 e_3)^* = -e_1\) であるから、ベクトルとバイベクトルの間に一対一対応が生まれる——3 次元空間特有の \(n=2k\) の自己双対性である。

\subsection{DST における双対写像}

第 II 部以降の時空代数 \(\mathrm{Cl}(3,1)\) では、4 次擬スカラー
\begin{equation}
  i = e_0 e_1 e_2 e_3
  \label{eq:dst-pseudoscalar}
\end{equation}
が DST の\textbf{双対写像} \(X \mapsto Xi\) の担い手となる。
この写像は通常時空セクターと双対時空セクターを入れ替える（第 IV 部で詳述）。

\(G(3,1,1)\) に拡張するとき、ヌルベクトル \(e_4\) は \(i\) と可換である：
\begin{equation}
  e_4 i = i e_4.
  \label{eq:e4-i-commute}
\end{equation}
したがって双対写像は \(\mathrm{Cl}(3,1)\) 部分代数上では従来どおり作用しつつ、\(e_4\) 方向は独立に保たれる。
これが PGA 版 DST において「双対性を保ちながら並進を内在化する」代数的基盤である。

5 次擬スカラー \(I_5 = e_0 e_1 e_2 e_3 e_4\) や、射影幾何代数文献で用いられる anti-scalar による双対化は、第 III 部以降の話題として残す。
本章で整えた反転・ノルム・階数分解は、次章の\textbf{ベルソル}と\textbf{ローター}——とくにサンドイッチ積 \(X' = RX\widetilde{R}\)——の代数的土台となる。

\section*{演習問題}

\begin{enumerate}[leftmargin=2em]
  \item \(M = 1 + e_1 + e_2 e_3 + 2e_1 e_2 e_3 \in G(3)\) を \(\langle M \rangle_k\)（\(k=0,1,2,3\)）に分解せよ。

  \item 直交単位ベクトル \(e_i,e_j\)（\(i\neq j\)）に対し、\(\widetilde{e_i e_j} = -e_i e_j\) および \(\widetilde{e_1 e_2 e_3} = e_1 e_2 e_3\) を \eqref{eq:reversion-sign} から確認せよ。

  \item \(G(3)\) において \(e_2^* = e_2 I_3^{-1}\) を計算し、結果が \(e_1 e_3\) となることを確かめよ。

  \item \(A_2 = e_1 e_2\)、\(B_2 = e_1 e_3\) の幾何積 \(A_2 B_2\) を計算し、\(\langle A_2 B_2 \rangle_0\) と \(\langle A_2 B_2 \rangle_2\) を求めよ。

  \item （発展）バイベクトル \(B\) が反転奇性 \(\widetilde{B}=-B\) を満たすとき、\(R = e^B\)（\(B^2\) がスカラー）が \(R\widetilde{R}=1\) を満たすことを示せよ。
  ヒント：\(e^B\) の級数展開で \(B\) の偶数次だけが寄与し、\(\widetilde{e^B} = e^{\widetilde{B}} = e^{-B}\) を用いよ。
\end{enumerate}

\chapter{ベルソルとローター}

前章では、多重ベクトルの階数分解、反転、ノルムを整えた。
本章では、これらの道具を使って\textbf{幾何変換を代数の元として書く}段階に進む。
可逆なベクトル積の全体を\textbf{ベルソル}（versor）と呼び、その中でも偶数階数部分に属する単位元を\textbf{ローター}（rotor）と呼ぶ。
ローターは連続的な回転（第 II 部ではローレンツ変換も）を表し、第 III 部以降の\textbf{モーター} \(M=\exp(\Omega)\) の原型となる。
DST 論文が前提とする \(R\widetilde{R}=1\) とサンドイッチ積 \(X'=MX\widetilde{M}\) も、ここで \(G(3)\) 上に確立する。

\section{ベルソルの定義と群構造}

\subsection{ベルソルの定義}

幾何代数 \(G(n)\) において、零でない可逆ベクトル \(v_1,\ldots,v_k\) の有限積
\begin{equation}
  V = v_1 v_2 \cdots v_k
  \label{eq:versor-def}
\end{equation}
を \textbf{ベルソル}と呼ぶ。
各 \(v_i\) は 1 次部分であり、\(V\) は一般に混合階数の多重ベクトルになる。
ベルソルは「ベクトルから生成される可逆元」であり、幾何変換の代数的数据を担う。

\paragraph{例}
\(G(3)\) において、単位ベクトル \(e_1,e_2\) の積 \(V = e_1 e_2\) は 2 次のベルソルである。
三つの単位ベクトルの積 \(W = e_1 e_2 e_3\) は 3 次のベルソルであり、反射を表す（次節）。

\subsection{逆元と群}

ベルソル \(V = v_1\cdots v_k\) の逆元は、因子の順序を逆にした積で与えられる：
\begin{equation}
  V^{-1} = v_k^{-1} \cdots v_1^{-1}.
  \label{eq:versor-inverse}
\end{equation}
ユークリッド設定では、ベクトル \(v\) の逆元は \eqref{eq:mv-inverse} より \(v^{-1} = \widetilde{v}/\|v\|^2\) である。
可逆ベクトルの有限積全体は、幾何積による乗法で\textbf{群}をなす——単位元は \(1\)、結合律は幾何積の結合律に従う。
この群は、直交群 \(O(n)\) の二重被覆群 \(\mathrm{Pin}(n)\) の代数的原型として機能する（群名は初出にとどめ、以降はベルソル・ローターという用語を主に用いる）。

\subsection{偶数ベルソルと奇数ベルソル}

ベルソル \(V = v_1\cdots v_k\) の\textbf{階数パリティ}（偶奇）は、因子の個数 \(k\) に依存する。
\(k\) が偶数なら \(V\) は\textbf{偶数ベルソル}——スカラーとバイベクトル（と高次元の偶数階成分）だけからなる——であり、連続的な回転の候補となる。
\(k\) が奇数なら \(V\) は\textbf{奇数ベルソル}——ベクトルや三ベクトルなど奇数階成分を含む——であり、一般に\textbf{反射}（向きの反転を伴う等長変換）を表す。

階数対合 \(\widehat{\cdot}\) \eqref{eq:grade-involution} との関係も明確である。
\(k\) 個のベクトルの積に対し
\begin{equation}
  \widehat{V} = \widehat{v_1}\cdots\widehat{v_k} = (-1)^k V,
  \label{eq:versor-grade-inv}
\end{equation}
すなわち偶数ベルソルは \(\widehat{V}=V\)、奇数ベルソルは \(\widehat{V}=-V\) となる。
第 2 章で強調した「偶数部分代数（四元数）と奇数部分の区別」は、ここで幾何変換の分類として再登場する。

\subsection{正規化と単位ベルソル}

ベルソル \(V\) のスカラー部ノルム \eqref{eq:mv-norm} が零でないとき、
\begin{equation}
  V_{\mathrm{u}} := \frac{V}{\sqrt{\|V\|^2}}
  \label{eq:versor-normalize}
\end{equation}
で\textbf{正規化}できる。
正規化されたベルソルは \(V\widetilde{V} = \pm 1\) を満たす。
符号は反射の回数（奇数ベルソルでは \(-1\)）に対応する。
\textbf{単位ベルソル}とは \(V\widetilde{V}=\pm 1\) かつ \(\|V\|^2=1\) を満たす正規化済みのベルソルを指す。

\section{ローターによる回転・反射の表現}

\subsection{反射}

単位ベクトル \(n\)（\(n^2=1\)、\(n^{-1}=\widetilde{n}=n\)）に対し、\(n\) に垂直な超平面での\textbf{反射}はベルソル \(n\) 自身で実現される：
\begin{equation}
  x' = -\,n\,x\,n.
  \label{eq:reflection}
\end{equation}
符号 \(-1\) は、反射が向きを反転する（行列式 \(-1\)）ことを反映している。
第 1 章演習で計算した \(e_1 e_2 e_1 = -e_2\) は、サンドイッチ \(n x n\)（先頭の \(-1\) を除いた形）に相当する。
完全な反射 \eqref{eq:reflection} では \(-e_1 e_2 e_1 = e_2\) となり、\(e_1\) に垂直な \(e_2\) 方向は固定される——法線方向 \(e_1\) 自身は \(e_1' = -e_1 e_1 e_1 = -e_1\) と反転する。

二つの反射の合成 \(V = n_2 n_1\) は偶数ベルソルとなり、一般に回転を与える。
\(G(3)\) では \(e_2 e_1 = -(e_1 e_2)\) が \(e_1 e_2\) 平面内の 180° 回転（同値に、平面法線方向の反射の二重適用）に対応する。

\subsection{ベルソル作用（サンドイッチ積）}

一般の可逆ベルソル \(V\) による\textbf{ベルソル作用}（versor action）は
\begin{equation}
  X' = V\,X\,V^{-1}
  \label{eq:versor-action}
\end{equation}
で定義される。
単位ベルソル \(V\widetilde{V}=\pm 1\) では \(V^{-1} = \pm\widetilde{V}\) であるから、符号に応じて
\begin{equation}
  X' = V\,X\,\widetilde{V}
  \quad\text{または}\quad
  X' = -V\,X\,\widetilde{V}
  \label{eq:sandwich}
\end{equation}
と書ける。
本書では主に\textbf{サンドイッチ積} \(X' = V X \widetilde{V}\) という形を用い、符号は \(V\widetilde{V}=\pm 1\) の取り方に吸収する。

この作用は次の性質を持つ。
\begin{itemize}[leftmargin=2em]
  \item \textbf{階数保存}：\(X\) が \(k\) 次なら \(X'\) も \(k\) 次である（ベルソルは偶数・奇数部分をそれぞれ偶数・奇数部分へ写す）。
  \item \textbf{ノルム保存}：\(\|X'\|^2 = \langle X'\widetilde{X'}\rangle_0 = \|X\|^2\)（単位ベルソルのとき）。
  \item \textbf{群作用}：\(V_2(V_1 X V_1^{-1})V_2^{-1} = (V_2 V_1)\,X\,(V_2 V_1)^{-1}\)——変換の合成はベルソルの積に対応する。
\end{itemize}
第 1 章の四元数による回転 \eqref{eq:quat-sandwich} は、\(G(3)\) の偶数単位ベルソル \(q\) に対する \eqref{eq:sandwich} の特殊例である。

\subsection{ローター}

\textbf{ローター} \(R\) とは、偶数の単位ベルソルで
\begin{equation}
  R\,\widetilde{R} = 1
  \label{eq:rotor-def}
\end{equation}
を満たす元である。
\eqref{eq:rotor-def} は第 2 章の \eqref{eq:rotor-unitarity} と同内容であり、回転（および第 II 部の真回転）が内積を保つ代数的根拠となる。

\(G(3)\) の偶数部分代数は四元数 \eqref{eq:quat-embedding} と同型であるから、第 1 章の単位四元数ローター \eqref{eq:quat-rotor} はすべて \(G(3)\) のローターである。
逆に、\(G(3)\) の任意のローターは、ある単位四元数 \(q\) として書ける。

\paragraph{例：\(e_1 e_2\) による平面内 90° 回転}
\(B = e_1 e_2\) は \(\widetilde{B} = -B\)、\(B^2 = -1\) を満たす。
ローター \(R = e^{-B\pi/4} = \cos(\pi/4) - B\sin(\pi/4) = \tfrac{1}{\sqrt{2}}(1 - e_1 e_2)\) に対し、
\(e_3\) 方向のベクトル \(e_3\) は \(B\) と可換であるから \(e_3' = R e_3 \widetilde{R} = e_3\)——\(e_1 e_2\) 平面に垂直な方向は固定される。
一方 \(e_1\) は \(e_1' = R e_1 \widetilde{R} = e_2\) となり、\(e_1 \to e_2\) の 90° 回転が実現される。

\subsection{回転の合成}

軸 \(\hat{u}_1\) 周りの角度 \(\theta_1\)、続いて軸 \(\hat{u}_2\) 周りの \(\theta_2\) の回転は、ローターの積
\begin{equation}
  R_{\mathrm{tot}} = R_2 R_1
  \label{eq:rotor-composition}
\end{equation}
で与えられる——四元数の非可換性 \eqref{eq:quat-mult} と同じ順序依存性である。
ジンバルロックに代表される特異点は、オイラー角のパラメータ化に固有のものであり、ローター \(R \in \mathrm{Spin}(3)\) 自体は滑らかな多様体上の点として扱える。

指数表示 \(R = e^{-\hat{u}\,\theta/2}\) \eqref{eq:quat-rotor} に現れる\textbf{半角} \(\theta/2\) は、ローター群 \(\mathrm{Spin}(3)\) が回転群 \(\mathrm{SO}(3)\) の二重被覆であることに由来する——\(2\pi\) 回転でローターの符号が反転し、\(4\pi\) で初めて \(1\) に戻る。

\paragraph{DST 論文との接続}
PGA 版 DST では、拡張モーター \(M = \exp(\Omega_{\mathrm{biv}}^{(5)})\) が \eqref{eq:rotor-def} と同型のユニタリ性 \(M\widetilde{M}=1\) を満たし、幾何変換は
\begin{equation}
  X' = M\,X\,\widetilde{M}
  \label{eq:motor-sandwich}
\end{equation}
で与えられる。
ねじれローター、ヌル並進、光円錐上の移動はすべてこのサンドイッチ積の枠内に収まる（第 III--V 部）。

\section{指数関数写像とリー代数的視点}

\subsection{バイベクトルによる指数写像}

バイベクトル \(B\) が \(B^2\) をスカラーとするとき、指数関数
\begin{equation}
  e^{B} = \sum_{n=0}^{\infty} \frac{B^n}{n!}
  \label{eq:exp-biv}
\end{equation}
は有限和に収束する（\(B^2 = s\) スカラーなら奇数・偶数項が分離される）。
\(B^2 = -\alpha^2 < 0\)（ユークリッド平面内の回転生成）のとき、単位方向 \(\hat{B} = B/\alpha\)（\(\hat{B}^2 = -1\)）を用いて
\begin{equation}
  e^{B} = \cos\alpha + \hat{B}\,\sin\alpha.
  \label{eq:exp-rotation}
\end{equation}
\(B^2 = \beta^2 > 0\)（第 II 部のミンコフスキー空間での\textbf{ブースト}生成）のときは
\begin{equation}
  e^{B} = \cosh\beta + \hat{B}\,\sinh\beta.
  \label{eq:exp-boost}
\end{equation}
となる。

\subsection{ローターの指数形}

平面内の回転角 \(\theta\) に対し、バイベクトル \(B = \tfrac{\theta}{2}\,\hat{u}\)（\(\hat{u}^2 = -1\)）を取ると
\begin{equation}
  R = e^{-B} = e^{-\hat{u}\,\theta/2}
  = \cos\frac{\theta}{2} - \hat{u}\,\sin\frac{\theta}{2}
  \label{eq:rotor-exp}
\end{equation}
となり、\eqref{eq:quat-rotor} と一致する。
半角が現れるのは、\eqref{eq:rotor-exp} が \(\mathrm{Spin}(3) \to \mathrm{SO}(3)\) の二重被覆写像であるためである。

\paragraph{反転との整合}
バイベクトルは反転奇性 \eqref{eq:biv-rev-odd} を満たすから、
\begin{equation}
  \widetilde{e^{B}} = e^{\widetilde{B}} = e^{-B},
  \label{eq:exp-reversion}
\end{equation}
したがって \(R = e^{-B}\) に対し
\begin{equation}
  R\,\widetilde{R} = e^{-B}\,e^{B} = 1.
  \label{eq:rotor-unitarity-proof}
\end{equation}
これは第 2 章演習（発展）の内容を本節で正式に確立したものである。
DST 論文の torsional rotor に対する \(R\widetilde{R}=1\) も、すべてのバイベクトル生成子が反転奇性 \((e_\mu e_\nu)^{\sim} = -e_\mu e_\nu\) を持つことに基づく——第 III 部でヌル生成子 \(N_\mu = e_4 \wedge e_\mu\) に拡張される。

\subsection{リー代数的視点}

ローター群の\textbf{リー代数}は、バイベクトル空間 \(\mathfrak{spin}(3) \cong \mathbb{R}^3\)（\(G(3)\) では \(\{e_1 e_2,\, e_2 e_3,\, e_3 e_1\}\) の張る 3 次元空間）である。
無限小ローター \(R = e^{\varepsilon B} \approx 1 + \varepsilon B\)（\(|B|\ll 1\)）による変換を 1 次で展開すると
\begin{equation}
  x' = e^{\varepsilon B}\,x\,e^{-\varepsilon B}
  \approx x + \varepsilon\,[B, x] + O(\varepsilon^2),
  \label{eq:infinitesimal-rotor}
\end{equation}
ここで \([B,x] := Bx - xB\) は\textbf{交換子}（commutator）である。
2 つのバイベクトル \(B_1,B_2\) の交換子 \([B_1,B_2]\) は再びバイベクトルであり、これが回転群の\textbf{合成}（非可換性）の無限小版を与える。

物理的には、\(B = \tfrac{1}{2}\sum_i \omega_i\, e_j e_k\)（\(\{i,j,k\}\) の順列）が角速度 \(\boldsymbol{\omega}\) に対応し、\eqref{eq:infinitesimal-rotor} は \(\dot{x} = \boldsymbol{\omega}\times x\) の代数版である。

\paragraph{第 II 部・第 III 部への先読み}
\(G(3,1)\) では 6 次元バイベクトル空間がローレンツ群のリー代数 \(\mathfrak{so}(3,1)\) を担い、\eqref{eq:exp-boost} の双曲関数がブーストを記述する。
\(G(3,1,1)\) ではヌルバイベクトル \(N_\mu\) が \(N_\mu^2 = 0\) かつ \(N_\mu N_\nu = 0\) を満たし、指数が
\begin{equation}
  e^{\sum_\mu \lambda^\mu N_\mu / 2}
  = 1 + \sum_\mu \frac{\lambda^\mu}{2}\,N_\mu
  \label{eq:exp-null-truncation}
\end{equation}
と\textbf{1 次で打ち切られる}——並進の有限生成と \(M\widetilde{M}=1\) の両方がここから生まれる（第 III 部・論文 Section 3）。

本章で確立した「ベルソル → ローター → 指数写像 → サンドイッチ積」の鎖は、第 II 部の \(G(3,1)\)（ローレンツ変換）と第 III 部の \(G(3,1,1)\)（モーター）へそのまま拡張される。

\section*{演習問題}

\begin{enumerate}[leftmargin=2em]
  \item \(G(3)\) において \(V = e_1 e_2\) が偶数ベルソルであることを、\(\widehat{V} = V\) \eqref{eq:versor-grade-inv} から確認せよ。
  また \(W = e_1 e_2 e_3\) が奇数ベルソルであり \(\widehat{W} = -W\) となることを示せよ。

  \item 単位ベクトル \(n = e_1\) に対し、反射 \eqref{eq:reflection} を \(x = e_1\) に適用し、\(x' = -e_1\) となることを直接計算で確かめよ。
  同じ \(n\) で \(x = e_2\) のとき \(x' = e_2\)（垂直方向は固定）となることも示せよ。

  \item ローター \(R = \tfrac{1}{\sqrt{2}}(1 - e_1 e_2)\) について、サンドイッチ積 \eqref{eq:sandwich} により \(e_1' = R e_1 \widetilde{R}\) を計算し、結果が \(e_2\) となることを示せよ。

  \item \(B = e_1 e_2\) に対し \eqref{eq:exp-rotation} を用いて \(R = e^{-B\pi/4}\) を求め、\eqref{eq:rotor-unitarity-proof} により \(R\widetilde{R}=1\) を確認せよ。

  \item バイベクトル \(B_1 = e_1 e_2\)、\(B_2 = e_2 e_3\) の交換子 \([B_1,B_2] = B_1 B_2 - B_2 B_1\) を計算せよ。
  結果が \(2\,e_1 e_3\) となり、再びバイベクトルであることを確かめよ。

  \item （発展）\(G(3,1)\) のバイベクトル \(H = e_0 e_1\)（\(e_0^2=-1\)、\(e_1^2=1\)）について \(H^2 = 1\) を示し、\eqref{eq:exp-boost} からブースト生成子 \(e^H\) の形を書け。
\end{enumerate}

\newpage

%==================================================
\part{\texorpdfstring{第II部　時空代数 \(G(3,1)\)}{第II部　時空代数 G(3,1)}}

\chapter{ミンコフスキー空間の幾何代数}

第 I 部ではユークリッド幾何代数 \(G(3)\) 上でベルソル・ローター・指数写像を整えた。
本章では符号 \((-,+,+,+)\) のミンコフスキー空間に対応する\textbf{時空代数} \(G(3,1)\)（\(\mathrm{Cl}(3,1)\) と同型）を導入する。
DST の親理論および PGA 版論文が前提とする四ベクトル基底、六つのバイベクトル生成子、ローレンツ変換のローター表現をここで確立する。
並進の内在化や \(G(3,1,1)\) への拡張は第 III 部に委ね、本章は \(\mathrm{Cl}(3,1)\) 部分代数の基礎に徹する。

\section{\texorpdfstring{符号 \((3,1)\) または \((1,3)\) の選択と物理的意味}{符号 (3,1) または (1,3) の選択と物理的意味}}

\subsection{ミンコフスキー符号とクリフォード代数}

4 次元ベクトル空間に符号 \((p,q,r,s)\) の二次形式を与えたとき、対応する幾何代数を \(G(p,q,r,s)\) と書く。
ミンコフスキー時空では 1 つの時間成分と 3 つの空間成分を持ち、符号は \((3,1)\) または \((1,3)\) のどちらかになる——前者は空間的ベクトル 3 個が \(+1\)、後者は時間的ベクトル 1 個が \(+1\) となる配置の違いである。

本書および DST 論文は、次の\textbf{ミンコフスキー符号} \((-,+,+,+)\) を一貫して採用する：
\begin{equation}
  e_0^2 = -1, \qquad e_1^2 = e_2^2 = e_3^2 = +1, \qquad \{e_\mu, e_\nu\} = 0 \quad (\mu \neq \nu).
  \label{eq:minkowski-signature}
\end{equation}
ここで \(e_0\) は時間方向、\(e_1,e_2,e_3\) は空間方向の直交単位ベクトルである。
反可換関係 \(\{e_\mu,e_\nu\}=0\) は、異なる方向の幾何積が外積のみを残すことを意味する（第 1 章 \eqref{eq:gp-decomp}）。

この代数を \textbf{\(G(3,1)\)} と呼ぶ。
文献によっては \(\mathrm{Cl}(3,1)\) とも書かれ、16 次元の\textbf{双四元数}代数と同型である（序章参照）。
\(G(3)\) が 8 次元だったのに対し、時間成分 \(e_0\) の追加により \(2^4 = 16\) 次元になる——多重ベクトルの階数分解は第 2 章と同様に
\[
G(3,1) = \langle 1 \rangle_0 \oplus \langle e_\mu \rangle_1 \oplus \langle e_\mu e_\nu \rangle_2 \oplus \langle e_\mu e_\nu e_\rho \rangle_3 \oplus \langle e_0 e_1 e_2 e_3 \rangle_4
\]
であり、各階数の次元は \(1,4,6,4,1\) である。

\subsection{\((3,1)\) と \((1,3)\) の違い}

符号 \((1,3)\) を採る文献では、\(e_0^2 = +1\)、\(e_i^2 = -1\)（\(i=1,2,3\)）と定義することがある。
これは \((3,1)\) との\textbf{符号の入れ替え}にすぎず、代数構造自体は同型である。
物理的内容——光円錐、ローレンツ不変量、因果律——は変わらない。

本書が \((-,+,+,+)\) を選ぶ理由は三つある。
第一に、特殊相対性理論の標準的記法 \(x^2 + y^2 + z^2 - (ct)^2\) との整合である。
第二に、DST 親理論および PGA 版論文が同一規約を用いること。
第三に、双四元数表現 \(e_0 = j\)（平方 \(-1\)）、\(e_i = kI,kJ,kK\)（平方 \(+1\)）との自然な対応である。

\paragraph{計量テンソル}
座標 \(x^\mu = (ct, x, y, z)\) に対し、ミンコフスキー計量は
\begin{equation}
  \eta_{\mu\nu} = \mathrm{diag}(-1, +1, +1, +1)
  \label{eq:minkowski-metric}
\end{equation}
である。
時空ベクトル \(X = x^\mu e_\mu\) の\textbf{二次形式}（スカラー部）は
\begin{equation}
  X^2 = \langle X^2 \rangle_0 = \eta_{\mu\nu}\, x^\mu x^\nu = -(ct)^2 + x^2 + y^2 + z^2.
  \label{eq:spacetime-norm}
\end{equation}
\(X^2 > 0\) は空間様、\(X^2 < 0\) は時間様、\(X^2 = 0\) は光的（null）ベクトルに対応する。

\subsection{四元数埋め込みとの連続性}

\(G(3)\) では空間バイベクトル \(I = e_3 e_2\)、\(J = e_1 e_3\)、\(K = e_2 e_1\) が四元数虚数単位 \eqref{eq:quat-embedding} として機能した。
\(G(3,1)\) では \(e_0\) が加わり、これらの \(I,J,K\) は\textbf{空間回転}生成子として残る。
一方 \(e_0\) を含むバイベクトル \(e_0 e_i\) は平方 \(+1\) を持ち、\textbf{ブースト}生成子となる——第 3 章 \eqref{eq:exp-boost} で予告した双曲関数の出現点である。

\section{時空ベクトルとバイベクトル}

\subsection{双四元数基底}

DST 論文および PGA 版論文は、四つのベクトル生成子を双四元数要素と同一視する：
\begin{equation}
  e_0 = j, \qquad e_1 = kI, \qquad e_2 = kJ, \qquad e_3 = kK.
  \label{eq:cl31-basis}
\end{equation}
\eqref{eq:minkowski-signature} が成り立つことは直接計算で確認できる。
例えば \(e_1 = kI\) より \(e_1^2 = kIkI = k^2 I^2 = (+1)(-1) = +1\) である。

一般の時空ベクトルは
\begin{equation}
  X = x^\mu e_\mu = ct\, e_0 + x\, e_1 + y\, e_2 + z\, e_3
  \label{eq:spacetime-vector}
\end{equation}
と書く。
\textbf{光的ベクトル}（null vector）は \eqref{eq:spacetime-norm} より
\begin{equation}
  L^2 = -(ct)^2 + x^2 + y^2 + z^2 = 0
  \label{eq:null-vector}
\end{equation}
を満たす。
光速で進む粒子の 4 速度はこの条件を満たす——第 III 部では、並進パラメータが光円錐上に制約される条件として再登場する。

\subsection{六つのバイベクトル生成子}

\(G(3,1)\) の 2 次部分は 6 次元であり、ローレンツ群のリー代数 \(\mathfrak{so}(3,1)\) を張る。
論文の分類に従い、六生成子は\textbf{双曲}（ブースト）と\textbf{巡回}（回転）の二群に分かれる。

\paragraph{双曲生成子（平方 \(+1\)）}
\begin{equation}
  iI = e_0 e_1, \qquad iJ = e_0 e_2, \qquad iK = e_0 e_3.
  \label{eq:hyperbolic-biv}
\end{equation}
各元は \(e_0\) と空間方向 \(e_i\) の外積であり、\(\mathfrak{so}(3,1)\) のブーストセクターを担う。
計算例：\(H = e_0 e_1\) に対し
\[
H^2 = e_0 e_1 e_0 e_1 = -e_0^2 e_1^2 = -(-1)(+1) = +1.
\]

\paragraph{巡回生成子（平方 \(-1\)）}
\begin{equation}
  I = e_3 e_2, \qquad J = e_1 e_3, \qquad K = e_2 e_1.
  \label{eq:cyclic-biv}
\end{equation}
これらは \(G(3)\) の四元数虚数単位 \eqref{eq:quat-embedding} と一致し、空間回転を生成する。
\(I^2 = e_3 e_2 e_3 e_2 = -e_3^2 e_2^2 = -1\) である。

表 \ref{tab:so31-generators} に六生成子の対応をまとめる。

\begin{table}[htbp]
  \centering
  \caption{\(G(3,1)\) の六バイベクトル生成子}
  \label{tab:so31-generators}
  \begin{tabular}{@{}llcl@{}}
    \toprule
    類 & 生成子 & 平方 & 物理的意味 \\
    \midrule
    双曲 & \(iI = e_0 e_1\), \(iJ = e_0 e_2\), \(iK = e_0 e_3\) & \(+1\) & ローレンツ・ブースト \\
    巡回 & \(I = e_3 e_2\), \(J = e_1 e_3\), \(K = e_2 e_1\) & \(-1\) & 空間回転 \\
    \bottomrule
  \end{tabular}
\end{table}

双曲生成子の平方が正、巡回生成子の平方が負——これは第 3 章 \eqref{eq:exp-rotation} と \eqref{eq:exp-boost} の分岐点であり、ミンコフスキー計量の符号 \((-,+,+,+)\) の直接的帰結である。

\subsection{擬スカラーと双対写像}

4 次擬スカラー
\begin{equation}
  i = e_0 e_1 e_2 e_3
  \label{eq:cl31-pseudoscalar}
\end{equation}
は \eqref{eq:dst-pseudoscalar} と同じ元である（\(i^2 = (e_0 e_1 e_2 e_3)^2 = -1\)）。
DST における\textbf{双対写像} \(X \mapsto Xi\) は、通常時空セクターと双対時空セクターを入れ替える——具体的には双曲生成子 \(\{iI,iJ,iK\}\) と巡回生成子 \(\{I,J,K\}\) を符号まで対応させる（第 IV 部で詳述）。

\(G(3,1,1)\) へ拡張するとき、ヌルベクトル \(e_4\) は \(i\) と可換である \eqref{eq:e4-i-commute}。
したがって双対写像は \(\mathrm{Cl}(3,1)\) 部分代数上では従来どおり作用し、並進方向 \(e_4\) は独立に保たれる——これが PGA 版 DST の代数的基盤である。

\subsection{電磁場のバイベクトル表現}

時空バイベクトルとして\textbf{電磁場} \(F\) を導入する。
空間成分 \(E_i\)（電場）と \(B_i\)（磁場）を用い
\begin{equation}
  F = E_1 e_0 e_1 + E_2 e_0 e_2 + E_3 e_0 e_3
    + B_1 e_2 e_3 + B_2 e_3 e_1 + B_3 e_1 e_2
  \label{eq:em-bivector}
\end{equation}
と書ける。
擬スカラー \(i\) を用いたコンパクトな表記
\begin{equation}
  F = \mathbf{E} + i\,\mathbf{B}
  \label{eq:em-compact}
\end{equation}
もよく使われる——ここで \(\mathbf{E} = E_i e_0 e_i\)、\(\mathbf{B} = B_i e_j e_k\)（循環順）である。
電場と磁場が単一の幾何対象 \(F\) に統合される点が、幾何代数的電磁気の出発点である。
マックスウェル方程式のコンパクトな形 \(\nabla F = J\) や相対論的変換則は次章で扱う。

\section{ローレンツ変換のローター表現}

\subsection{ローターとサンドイッチ積}

第 3 章で確立したローターの定義 \eqref{eq:rotor-def} とサンドイッチ積 \eqref{eq:sandwich} は、\(G(3,1)\) 上でもそのまま有効である。
偶数単位ベルソル \(R\) が \(R\widetilde{R} = 1\) を満たすとき、時空ベクトル \(X\) の変換は
\begin{equation}
  X' = R\,X\,\widetilde{R}
  \label{eq:lorentz-sandwich}
\end{equation}
で与えられる。
この作用は階数 1 を保存し、\eqref{eq:spacetime-norm} の二次形式を不変にする——すなわちローレンツ変換である。

\subsection{空間回転}

空間 \(e_i e_j\) 平面内の回転は、第 3 章と同様に \(B^2 = -\alpha^2 < 0\) なるバイベクトル \(B\) から
\begin{equation}
  R = e^{-B} = \cos\alpha + \hat{B}\sin\alpha
  \label{eq:rotation-rotor}
\end{equation}
（\eqref{eq:exp-rotation}）で与えられる。
例えば \(e_1 e_2\) 平面内の角度 \(\theta\) の回転は \(B = \tfrac{\theta}{2} e_1 e_2\) より \(R = \cos(\theta/2) - e_1 e_2 \sin(\theta/2)\) である。

\subsection{ローレンツ・ブースト}

双曲バイベクトル \(H = e_0 e_1\) は \(H^2 = +1\) であるから、第 3 章 \eqref{eq:exp-boost} より
\begin{equation}
  R_{\mathrm{bst}} = e^{-H\phi/2} = \cosh\frac{\phi}{2} - H\,\sinh\frac{\phi}{2}
  \label{eq:boost-rotor}
\end{equation}
がブースト生成ローターを与える。
ここで \(\phi\) は\textbf{ラピディティ}（rapidity）であり、速度 \(v\) との関係は \(\phi = \operatorname{artanh}(v/c)\) である。
\(H\) の符号と向きがブースト方向を決める。

\paragraph{例：\(x\) 方向へのブースト}
\(H = e_0 e_1\)、\(\widetilde{H} = -H\) として \eqref{eq:boost-rotor} より
\[
R = \cosh\frac{\phi}{2} - e_0 e_1 \sinh\frac{\phi}{2}.
\]
サンドイッチ積 \eqref{eq:lorentz-sandwich} を \(X = ct\, e_0 + x\, e_1\) に適用すると
\begin{align}
  (ct)' &= ct\,\cosh\phi - x\,\sinh\phi, \\
  x'    &= -ct\,\sinh\phi + x\,\cosh\phi,
\end{align}
すなわち標準的なローレンツ・ブーストが得られる。

\subsection{回転とブーストの統一}

回転（\(B^2 < 0\)、三角関数）とブースト（\(H^2 > 0\)、双曲関数）が、同一の指数形式 \(R = e^{-B/2}\) で記述される——これが幾何代数的特殊相対性理論の核心である。
一般のローレンツ変換は、適切な 6 次元バイベクトル \(\Omega \in \mathfrak{so}(3,1)\) から
\begin{equation}
  R = e^{-\Omega/2}, \qquad X' = R\,X\,\widetilde{R}
  \label{eq:general-lorentz}
\end{equation}
として与えられる。

\subsection{無限小変換とリー代数}

\(|B| \ll 1\) のとき \(R = e^{-B/2} \approx 1 - B/2\) より、無限小ローレンツ変換は
\begin{equation}
  X' \approx X - \tfrac{1}{2}[B, X] + O(B^2)
  \label{eq:infinitesimal-lorentz}
\end{equation}
（\eqref{eq:infinitesimal-rotor} の \(G(3,1)\) 版）で与えられる。
六つのバイベクトル生成子 \eqref{eq:hyperbolic-biv}–\eqref{eq:cyclic-biv} の交換子は再び \(\mathfrak{so}(3,1)\) 内に閉じ、ポアンカレ群の\textbf{非並進}部分 \(\mathrm{SO}(3,1)\) のリー代数構造を与える。

\subsection{DST 論文との接続}

PGA 版 DST では、\(\mathrm{Cl}(3,1)\) 部分代数の六バイベクトルが二つの役割を担う。
\textbf{双曲生成子} \(\{iI,iJ,iK\}\) は通常時空側のブースト——ねじれ不整合の「引っ張る」方向——に対応する。
\textbf{巡回生成子} \(\{I,J,K\}\) は双対時空側の回転——コンパクトな位相情報の保存——に対応する。
擬スカラー \(i\) による双対写像 \(X \mapsto Xi\) が両セクターを入れ替える（第 IV 部）。

第 III 部では、ヌルベクトル \(e_4\) を加えた \(G(3,1,1)\) 上で四つの\textbf{ヌル生成子} \(N_\mu = e_4 \wedge e_\mu\) が並進を内在化する。
六生成子に四生成子を加えた 10 次元バイベクトル空間がポアンカレ代数 \(\mathfrak{iso}(3,1)\) となり、回転・ブースト・並進が単一のモーター \(M = \exp(\Omega)\) に統合される——本章で整えた \(G(3,1)\) のローター理論が、その直接的な土台となる。

\section*{演習問題}

\begin{enumerate}[leftmargin=2em]
  \item \(H = e_0 e_1\) について \(H^2 = +1\) を \eqref{eq:minkowski-signature} から直接計算で示せ。
  同様に \(B = e_1 e_2\) について \(B^2 = -1\) を確認せよ。

  \item 光的ベクトル \(L = c t_0 e_0 + x_0 e_1\)（\(x_0 \neq 0\)）について \eqref{eq:null-vector} より \(L^2 = 0\) となる \(t_0\) の値を求めよ。
  結果が \(t_0 = x_0/c\) となり、光速条件と一致することを確認せよ。

  \item ブーストローター \eqref{eq:boost-rotor} について、\(H\) の反転奇性 \(\widetilde{H} = -H\) から \(\widetilde{R} = \cosh(\phi/2) + H\sinh(\phi/2)\) を求め、\(R\widetilde{R} = 1\) を示せ。

  \item \(R = \cosh(\phi/2) - e_0 e_1 \sinh(\phi/2)\) について、サンドイッチ積 \eqref{eq:lorentz-sandwich} により \(e_0' = R e_0 \widetilde{R}\) と \(e_1' = R e_1 \widetilde{R}\) を計算し、\(e_0\) と \(e_1\) が双曲関数で混合されることを示せ。

  \item 双曲生成子 \(iI = e_0 e_1\) と巡回生成子 \(I = e_3 e_2\) の交換子 \([iI, I]\) を計算せよ。
  結果が再びバイベクトルであることを確認し、\(\mathfrak{so}(3,1)\) の閉性を確かめよ。

  \item （発展）擬スカラー \(i = e_0 e_1 e_2 e_3\) について \(iI = e_0 e_1\) との積 \(i \cdot (iI)\) を計算し、巡回生成子 \(I = e_3 e_2\) との関係を調べよ。
  双対写像 \(X \mapsto Xi\) が双曲生成子と巡回生成子を入れ替えることの手がかりとなる。
\end{enumerate}

\chapter{時空代数の応用}

前章で \(G(3,1)\) 上のローター、電磁場バイベクトル \(F=\mathbf{E}+i\mathbf{B}\)、サンドイッチ積 \eqref{eq:lorentz-sandwich} を整えた。
本章では、これらを\textbf{相対論的運動学}、\textbf{マックスウェル方程式}、\textbf{ディラック方程式}へ接続する——特殊相対性理論と古典電磁気論、そして量子力学との橋渡しである。
DST 固有の双対時空や \(G(3,1,1)\) への拡張は第 III--IV 部に委ね、ここでは \(\mathrm{Cl}(3,1)\) 部分代数上の標準的 STA 応用に徹する。

\section{相対論的運動学の書き換え}

\subsection{世界線と固有時間}

質量 \(m>0\) の粒子の\textbf{世界線}は、時空ベクトル \eqref{eq:spacetime-vector} の軌道 \(X(\tau)\) として与えられる。
パラメータ \(\tau\) は\textbf{固有時間}（proper time）であり、粒子が携える時計が刻む時間に対応する。
固有時間は座標時間 \(t\) とは一般に異なる——静止系では \(\tau = t\) だが、運動系では \(\tau < t\) となる。

世界線上の任意の点で、接ベクトル \(v = dX/d\tau\) は\textbf{時間様}（timelike）——\(v^2 < 0\)——である。
ローター \eqref{eq:lorentz-sandwich} は \(X^2\) を不変に保つ。

\subsection{四速度と正規化}

\textbf{四速度}（4-velocity）は世界線の接ベクトル
\begin{equation}
  v = \frac{dX}{d\tau} = c\,\frac{dt}{d\tau}\,e_0 + \frac{dx}{d\tau}\,e_1 + \frac{dy}{d\tau}\,e_2 + \frac{dz}{d\tau}\,e_3
  \label{eq:four-velocity}
\end{equation}
で定義される。
\(v\) は階数 1 の多重ベクトル——すなわち普通の時空ベクトル——である。

粒子の静止系では \(v = c\,e_0\) となる。
一般の慣性系では、空間速度 \(\boldsymbol{\beta} = \mathbf{v}_{\mathrm{sp}}/c\) を用いて
\begin{equation}
  v = \gamma c\, e_0 + \gamma c\,\beta_x e_1 + \gamma c\,\beta_y e_2 + \gamma c\,\beta_z e_3,
  \qquad \gamma = \frac{1}{\sqrt{1-\beta^2}}
  \label{eq:four-velocity-gamma}
\end{equation}
と書ける。
四速度の\textbf{正規化}は
\begin{equation}
  v^2 = \langle v^2 \rangle_0 = -c^2
  \label{eq:four-velocity-norm}
\end{equation}
である——静止系の \(v=c e_0\) で \(v^2 = c^2 e_0^2 = -c^2\) となり、\eqref{eq:four-velocity-gamma} でも同値に確認できる。

\paragraph{固有時間と座標時間}
座標時間 \(t\) との関係は \(d\tau = dt/\gamma\) である。
\(v^2=-c^2\) を成分表示 \((\gamma c, \gamma c\boldsymbol{\beta})\) に代入すると \(\gamma^2(1-\beta^2)=1\) が再現され、ローレンツ因子 \(\gamma\) の幾何学的意味が明確になる。

\subsection{四運動量}

\textbf{四運動量}は
\begin{equation}
  p = m v
  \label{eq:four-momentum}
\end{equation}
で定義される。
\(m\) は静止質量、\(v\) は \eqref{eq:four-velocity-norm} を満たす四速度である。
直ちに
\begin{equation}
  p^2 = m^2 v^2 = -m^2 c^2
  \label{eq:four-momentum-norm}
\end{equation}
が従い、これが相対論的エネルギー–運動量関係の幾何代数版である。

\paragraph{エネルギーと運動量の読み取り}
\(x\) 方向に速度 \(\beta c\) で運動する粒子 \eqref{eq:four-velocity-gamma} では
\begin{equation}
  p = \gamma m c\, e_0 + \gamma m \beta c\, e_1.
  \label{eq:momentum-example}
\end{equation}
スカラー部 \(\langle p e_0 \rangle_0 = -\gamma m c\) より、エネルギーは
\begin{equation}
  E = -\langle p e_0 \rangle_0\, c = \gamma m c^2,
  \label{eq:energy-ga}
\end{equation}
空間運動量は \(p_x = \langle p e_1 \rangle_0 = \gamma m \beta c\) である——テンソル記法の \(p^\mu = (E/c,\, p_x,\, p_y,\, p_z)\) と一致する。

\subsection{加速度と直交性}

四速度の \(\tau\) 微分を\textbf{四加速度}と呼ぶ：
\begin{equation}
  a = \frac{dv}{d\tau}.
  \label{eq:four-acceleration}
\end{equation}
\(v^2 = -c^2\) が \(\tau\) に依存しないことから、
\begin{equation}
  0 = \frac{d(v^2)}{d\tau} = v a + a v = 2\, v\cdot a,
  \label{eq:velocity-acceleration-ortho}
\end{equation}
すなわち \(v\cdot a = 0\) である。
四速度と四加速度は、ミンコフスキー内積の意味で常に直交する——これは「速度を変える方向は、現在の速度方向に垂直である」という直感の相対論版である。

\subsection{ローレンツ変換下の四速度}

四速度は階数 1 のベクトルであるから、座標変換はサンドイッチ積 \eqref{eq:lorentz-sandwich} で与えられる：
\begin{equation}
  v' = R\,v\,\widetilde{R}.
  \label{eq:four-velocity-transform}
\end{equation}
前章の \(x\) 方向ブースト \eqref{eq:boost-rotor} を \(v = c e_0\)（静止四速度）に適用すると
\begin{equation}
  v' = \gamma c\, e_0 + \gamma \beta c\, e_1,
  \label{eq:boost-four-velocity}
\end{equation}
すなわち \eqref{eq:four-velocity-gamma} が得られる。
回転とブーストが同一のローター形式 \(R=e^{-\Omega/2}\) \eqref{eq:general-lorentz} で記述される点が、幾何代数的 SRT の核心である。

\subsection{電磁場中の運動}

電荷 \(q\) の粒子が電磁場 \eqref{eq:em-compact} 中を運動するとき、運動方程式は
\begin{equation}
  m\,\frac{dv}{d\tau} = q\, F\cdot v
  \label{eq:lorentz-force}
\end{equation}
で与えられる。
右辺の \(F\cdot v\) はバイベクトル \(F\) とベクトル \(v\) の\textbf{内積} \eqref{eq:gp-decomp} により階数 1 のベクトルを与える——これが\textbf{ローレンツ力}の STA 形である。

\paragraph{例：一様磁場中の円運動}
\(z\) 方向一様磁場 \(\mathbf{B} = B e_1 e_2\)（\(B>0\)）、初期四速度 \(v = \gamma c e_0 + \gamma v_0 e_1\) を考える。
電場が零なら \(F = i\mathbf{B}\) である。
\(F\cdot v\) を計算すると、\(v\) の空間成分が \(e_2\) 方向へ曲がる項が現れ、円運動のサイクロトロン半径 \(r = \gamma m v_0 / (qB)\) が得られる——成分計算の詳細は演習に委ねる。

\paragraph{DST への先読み}
質量零粒子の四速度は光的ベクトル \eqref{eq:null-vector} 上に制約される。
並進自由度は \(\mathrm{Cl}(3,1)\) だけでは外付け的だったが、第 III 部では \(G(3,1,1)\) のモーターとして内在化される。

\section{電磁気のマックスウェル方程式のコンパクトな形}

\subsection{時空勾配と四ベクトル電流}

座標 \(x^\mu = (ct, x, y, z)\) に対し、\textbf{時空勾配}（spacetime gradient）は
\begin{equation}
  \nabla = e^\mu \partial_\mu
  = e_0 \frac{\partial}{\partial(ct)} + e_1 \frac{\partial}{\partial x} + e_2 \frac{\partial}{\partial y} + e_3 \frac{\partial}{\partial z}
  \label{eq:spacetime-gradient}
\end{equation}
で定義される。
ここで \(e^\mu\) は \(e_\mu\) の逆元（\(e^0 = e_0\)、\(e^i = e_i\)）である。

\textbf{四ベクトル電流}は
\begin{equation}
  J = \rho\, e_0 + j_x e_1 + j_y e_2 + j_z e_3
  \label{eq:four-current}
\end{equation}
と書く——\(\rho\) は電荷密度、\(\mathbf{j} = j_x e_1 + j_y e_2 + j_z e_3\) は電流密度である。

\subsection{マックスウェル方程式の統合形}

前章 \eqref{eq:em-compact} で導入した電磁場バイベクトル \(F\) に対し、マックスウェル方程式は\textbf{単一方程式}
\begin{equation}
  \nabla F = J
  \label{eq:maxwell-compact}
\end{equation}
として書ける。
左辺は幾何積——勾配（ベクトル）と場（バイベクトル）の積——であり、結果は一般に混合階数の多重ベクトルになる。

\paragraph{階数分解}
\eqref{eq:maxwell-compact} を階数ごとに分離すると、二つの物理的内容が現れる：
\begin{align}
  \nabla\cdot F &= J, \label{eq:maxwell-vector} \\
  \nabla\wedge F &= 0. \label{eq:maxwell-bivector}
\end{align}
\(\nabla\cdot F\) は\textbf{内積}（\(\nabla\) と \(F\) のベクトル–バイベクトル contraction）で階数 1 を与え、\(\nabla\wedge F\) は\textbf{外積}で階数 2 を与える。

\eqref{eq:maxwell-vector} はガウスの法則 \(\nabla\cdot\mathbf{E} = \rho/\varepsilon_0\) とアンペールの法則 \(\nabla\times\mathbf{B} = \mu_0\mathbf{j} + \partial\mathbf{E}/\partial t\) を統合する（自然単位系 \(\varepsilon_0=\mu_0=1\) では \(\rho\) と \(\mathbf{j}\) がそのまま現れる）。
\eqref{eq:maxwell-bivector} はファラデーの法則 \(\nabla\times\mathbf{E} = -\partial\mathbf{B}/\partial t\) と磁荷なし条件 \(\nabla\cdot\mathbf{B}=0\) を統合する。
四つのマックスウェル方程式が、ベクトル部とバイベクトル部の二つに圧縮される——これが STA 記法の利点である。

\subsection{ローレンツ不変量}

電磁場の\textbf{ローレンツ不変量}は \(F\) の幾何積 \(F^2\) のスカラー部と擬スカラー部に分解される：
\begin{equation}
  F^2 = \langle F^2 \rangle_0 + i\,\langle F^2 \rangle_4
  = (\mathbf{E}^2 - \mathbf{B}^2) + 2i\,\mathbf{E}\cdot\mathbf{B}.
  \label{eq:em-invariants}
\end{equation}
\(\mathbf{E}^2 - \mathbf{B}^2\) は不変量 \(F_{\mu\nu}F^{\mu\nu}\) に、\(\mathbf{E}\cdot\mathbf{B}\) は \(F_{\mu\nu}\tilde{F}^{\mu\nu}\) に対応する。
ローレンツ変換 \(F' = RF\widetilde{R}\) \eqref{eq:lorentz-sandwich} の下で \(F^2\) は不変——すなわち \(F^2 = F'^2\) である。

\paragraph{例：静止系から見た純粋電場}
静止系で \(\mathbf{E} = E e_0 e_1\)、\(\mathbf{B}=0\) のとき \(F^2 = E^2 > 0\)（空間様）。
\(x\) 方向ブースト \eqref{eq:boost-rotor} 後の系では \(F' = RF\widetilde{R}\) により電場と磁場が混合され、\(F'^2 = E^2\) が保存される——\(E'\) と \(B'\) の具体的な値は演習で求める。

\subsection{ポテンシャルとゲージ}

電磁場は一般に\textbf{ポテンシャル} \(A\)（1 次の時空ベクトル）から
\begin{equation}
  F = \nabla\wedge A
  \label{eq:em-potential}
\end{equation}
と書ける。
\(\nabla\wedge F = \nabla\wedge(\nabla\wedge A) = 0\) は恒等式 \eqref{eq:maxwell-bivector} を自動的に満たす。

ゲージ変換 \(A \mapsto A + \nabla\chi\) は \(F\) を不変にする。
\textbf{Lorenz ゲージ} \(\nabla\cdot A = 0\) の下では、波動方程式 \(\nabla^2 A = J\) が従い、電磁波の伝播が記述できる——詳細はここでは省略する。

\section{ディラックスピノルとの接点}

\subsection{双四元数とスピノル空間}

第 I 章で述べたように、\(G(3,1) \cong \mathrm{Cl}(3,1)\) は\textbf{双四元数} \(\mathbb{H}\oplus\mathbb{H}\) と同型である。
16 次元の実代数は、4 次元の\textbf{複素}ベクトル空間——すなわちディラックスピノル空間——として解釈できる。

従来の量子力学では、スピノルを「代数の上に載せた外部の列ベクトル」として導入し、ガンマ行列 \(\gamma^\mu\) でローレンツ共変性を課す。
STA の視点では、スピノルは\textbf{代数自身の元}として現れる——具体には、\textbf{偶数部分代数}の元、または\textbf{最小左イデアル}の元として記述する。

\subsection{ガンマ行列と基底ベクトルの対応}

本書の双四元数基底 \eqref{eq:cl31-basis} において、ディラックの \(\gamma\) 行列は四ベクトル生成子と同一視できる：
\begin{equation}
  \gamma^0 \longleftrightarrow e_0 = j, \qquad
  \gamma^1 \longleftrightarrow e_1 = kI, \qquad
  \gamma^2 \longleftrightarrow e_2 = kJ, \qquad
  \gamma^3 \longleftrightarrow e_3 = kK.
  \label{eq:gamma-correspondence}
\end{equation}
クリフォード関係 \(\{e_\mu, e_\nu\} = 2\eta_{\mu\nu}\) \eqref{eq:minkowski-signature} は、標準的な \(\{\gamma^\mu, \gamma^\nu\} = 2\eta^{\mu\nu}\) と一致する。

\subsection{ディラック演算子}

時空勾配 \eqref{eq:spacetime-gradient} から\textbf{ディラック演算子}（slash 演算子）を
\begin{equation}
  \not\!\partial = \nabla = e^\mu \partial_\mu
  \label{eq:dirac-operator}
\end{equation}
と定義する。
自由粒子の\textbf{ディラック方程式}は
\begin{equation}
  (\not\!\partial + m)\,\psi = 0
  \label{eq:dirac-equation}
\end{equation}
である——\(\psi\) はスピノル、\(m\) は静止質量である。

\(\not\!\partial\) の二乗は
\begin{equation}
  \not\!\partial^2 = \nabla^2 = \partial_\mu \partial^\mu
  \label{eq:dirac-squared}
\end{equation}
となり、\eqref{eq:dirac-equation} からクライン–ゴルドン方程式 \((\nabla^2 - m^2)\psi = 0\) が従う。

\subsection{ローターによるスピノル変換}

時空ベクトル \(X\) のローレンツ変換はサンドイッチ積 \(X' = RX\widetilde{R}\) \eqref{eq:lorentz-sandwich} であるが、スピノル \(\psi\) への作用は\textbf{左乗算}
\begin{equation}
  \psi' = R\,\psi
  \label{eq:spinor-transform}
\end{equation}
で与えられる——サンドイッチ積ではない点に注意する。
これはスピノルが \(\mathrm{Spin}(3,1)\) の\textbf{基本表現}（左イデアル上の表現）に属するためである。
\eqref{eq:spinor-transform} の下で \eqref{eq:dirac-equation} は共変——すなわち \(\psi\) を変換した後も同じ形の方程式が成り立つ。

偶数ローター \(R\) は \(\mathrm{Spin}(3,1)\) の元であり、\(R\widetilde{R}=1\) \eqref{eq:rotor-def} を満たす。
回転 \(R = e^{-B/2}\) \eqref{eq:rotation-rotor} もブースト \(R_{\mathrm{bst}}\) \eqref{eq:boost-rotor} も、同一の \eqref{eq:spinor-transform} でスピノルに作用する。

\subsection{最小左イデアル（概観）}

より具体的には、原始冪等元
\begin{equation}
  P = \frac{1 + e_0}{2}\cdot\frac{1 + i}{2}
  \label{eq:spinor-idempotent}
\end{equation}
（\(i = e_0 e_1 e_2 e_3\) \eqref{eq:cl31-pseudoscalar}）が\textbf{最小左イデアル}
\begin{equation}
  \mathcal{S} = G(3,1)\,P
  \label{eq:minimal-left-ideal}
\end{equation}
を生成する。
\(\mathcal{S}\) の一般元 \(\psi\) は 4 つの実パラメータ——したがって 4 つの複素成分——を持ち、標準的ディラックスピノルと同次元である。

擬スカラー \(i\) を量子力学の虚数単位と同一視すると、\(\psi\) はウェイル基底の 4 成分複素列として読める。
キラル射影 \(P_L = (1-i)/2\)、\(P_R = (1+i)/2\) も代数内部から自然に定義される——詳細は第 IV 部「スピノル・量子状態との統合」で展開する。

\paragraph{第 IV 部への予告}
DST では、通常時空ローター \(R_{\mathrm{usual}}\) と双対時空ローター \(R_{\mathrm{dual}}\) の積 \(R_{\mathrm{total}} = R_{\mathrm{usual}} R_{\mathrm{dual}}\) がスピノル \eqref{eq:spinor-transform} に作用する。
質量項 \(m\psi\) は二つのローターの\textbf{ねじれ不整合}から幾何学的に生じ、\(\bar\psi\psi\) はその期待値に比例する——これらは第 IV--V 部の中心テーマである。
本章で確立した \(\gamma^\mu \leftrightarrow e_\mu\) の対応と \eqref{eq:dirac-equation} が、その代数的土台となる。

\section*{演習問題}

\begin{enumerate}[leftmargin=2em]
  \item \(v = \gamma c e_0 + \gamma \beta c e_1\) \eqref{eq:four-velocity-gamma} について、\(v^2 = -c^2\) \eqref{eq:four-velocity-norm} を直接計算で確認せよ。
  同様に \(v\cdot a = 0\) \eqref{eq:velocity-acceleration-ortho} が \(a = dv/d\tau\) から従うことを示せ。

  \item 四運動量 \eqref{eq:momentum-example} について、\(p^2 = -m^2 c^2\) \eqref{eq:four-momentum-norm} を成分計算で確認せよ。
  エネルギー \(E = \gamma mc^2\) \eqref{eq:energy-ga} と運動量 \(p_x = \gamma m \beta c\) が得られることも示せ。

  \item \(z\) 方向一様磁場 \(F = B e_1 e_2\)（\(B>0\)）中で、\(v = c e_0\) から出発した粒子に \eqref{eq:lorentz-force} を適用せよ。
  \(F\cdot v\) を計算し、加速度が \(e_2\) 方向を向くことを確認せよ。

  \item 静止系で \(F = E e_0 e_1\)（\(\mathbf{B}=0\)）のとき \(F^2 = E^2\) \eqref{eq:em-invariants} を示せ。
  ブーストローター \(R = \cosh(\phi/2) - e_0 e_1 \sinh(\phi/2)\) \eqref{eq:boost-rotor} により \(F' = RF\widetilde{R}\) を計算し、\(F'^2 = E^2\) が保存されることを確認せよ。

  \item \(\nabla\wedge(\nabla\wedge A) = 0\) が任意の \(A\) について恒真であることを、\(\nabla\wedge\) の定義 \eqref{eq:gp-decomp} から示せ。
  これが \eqref{eq:maxwell-bivector} を自動的に満たす理由を説明せよ。

  \item 基底 \eqref{eq:gamma-correspondence} について \(\{e_0, e_1\} = 2 e_0 e_1\)（反可換）と \(\{e_1, e_2\} = 0\)（可換）を \eqref{eq:minkowski-signature} から確認せよ。
  \(\not\!\partial^2 = \nabla^2\) \eqref{eq:dirac-squared} の意味を、成分表示で説明せよ。

  \item （発展）偶数ローター \(R = \cos(\theta/2) - e_1 e_2 \sin(\theta/2)\) \eqref{eq:rotation-rotor} がスピノル \(\psi = 1 + P\)（\(P\) \eqref{eq:spinor-idempotent}）に \eqref{eq:spinor-transform} で作用するとき、\(\psi' = R\psi\) を計算せよ。
  ベクトル変換 \(X' = RX\widetilde{R}\) \eqref{eq:lorentz-sandwich} との違いを述べよ。
\end{enumerate}

\newpage

%==================================================
\part{\texorpdfstring{第III部　射影幾何代数 \(G(3,1,1)\)}{第III部　射影幾何代数 G(3,1,1)}}

\chapter{ヌル次元の導入と射影構造}

第 II 部で整えた時空代数 \(G(3,1)\) は、ローレンツ変換をローターとして表現する——しかし並進は代数の内部からは生まれない。
本章では、平方がゼロの\textbf{ヌルベクトル} \(e_4\) を一つ付加して射影幾何代数 \(G(3,1,1)\) を構築し、点・直線・平面の符号化と join/meet 操作を整える。
並進生成子 \(N_\mu = e_4 \wedge e_\mu\) の代数的性質とモーター構成は次章に委ねる。
本章の役割は、PGA 版 DST が前提とする\textbf{射影構造}の土台を置くことにある。

\section{\texorpdfstring{\(G(3,1)\) にヌルベクトル \(e_4\) を加える}{G(3,1) にヌルベクトル e4 を加える}}

\subsection{動機：並進の内在化}

\(\mathrm{Cl}(3,1)\) 上では、回転とブースト——すなわち \(\mathfrak{so}(3,1)\) の六自由度——は六つのバイベクトル生成子 \eqref{eq:hyperbolic-biv}・\eqref{eq:cyclic-biv} として自然に現れる。
一方、時空並進 \(x^\mu \mapsto x^\mu + a^\mu\) は、座標変換として外付け的に与えられていた。
PGA 版 DST では、この並進を代数の内部に内在化する必要がある。

射影幾何代数の標準的作法は、既存のベクトル空間に\textbf{退化ベクトル}（degenerate vector）——平方がゼロのベクトル——を一つ加えることである。
Selig 以来の剛体運動学や Lengyel--Gunn 系 PGA 文献でも、並進はヌルベクトルと既存ベクトルの外積からなるバイベクトルとして符号化される。
本書では、ミンコフスキー時空の四ベクトル \eqref{eq:cl31-basis} に、新しいベクトル \(e_4\) を付加して \(G(3,1,1)\) へ持ち上げる。

\subsection{ヌルベクトル \(e_4\) の定義}

\(e_4\) は次の関係で特徴づけられる：
\begin{equation}
  e_4^2 = 0, \qquad \{e_4, e_\mu\} = 0 \quad (\mu = 0,1,2,3).
  \label{eq:e4-definitions}
\end{equation}
第一式は \(e_4\) が\textbf{ヌル}（null）——自己内積がゼロ——であることを意味する。
第二式は \(e_4\) が既存の四ベクトル \(e_0,e_1,e_2,e_3\) とすべて反可換することを意味する。
したがって \(e_4\) は \(\mathrm{Cl}(3,1)\) 部分代数の生成元ではなく、独立な\textbf{射影方向}を担う。

\paragraph{計算例}
\(e_4\) と \(e_1 = kI\) の積を考える。
反可換関係より \(e_4 e_1 = -e_1 e_4\) である。
さらに \(e_4^2 = 0\) から \(e_4 e_1 e_4 = -e_4 e_1 e_4\) となり、\(2\,e_4 e_1 e_4 = 0\) すなわち \(e_4 e_1 e_4 = 0\) が従う——この「サンドイッチ消滅」は、後に現れるヌルバイベクトルの代数的性質の原型となる。

\subsection{\(G(3,1,1)\) の次元構造}

5 つの直交生成元 \(\{e_0,e_1,e_2,e_3,e_4\}\) から、クリフォード代数 \(G(3,1,1)\) が生成される。
ベクトル空間は 5 次元、代数全体は \(2^5 = 32\) 次元である。
\(\{e_0,e_1,e_2,e_3\}\) だけから生成される部分代数は 16 次元の \(\mathrm{Cl}(3,1)\)——双四元数代数——に一致する。

階数分解は次のとおり：
\begin{equation}
  G(3,1,1)
  = \langle 1 \rangle_0 \oplus \langle e_a \rangle_1 \oplus \langle e_a e_b \rangle_2
  \oplus \langle e_a e_b e_c \rangle_3 \oplus \langle e_a e_b e_c e_d \rangle_4
  \oplus \langle I_5 \rangle_5,
  \label{eq:g311-grade-decomp}
\end{equation}
ここで添字 \(a,b,c,d\) は \(\{0,1,2,3,4\}\) から選ぶ。
各階数の次元は
\begin{equation}
  1,\ 5,\ 10,\ 10,\ 5,\ 1
  \label{eq:g311-grade-dims}
\end{equation}
である。
とくに 2 次部分は \(\binom{5}{2} = 10\) 次元——第 II 部の六生成子に、\(e_4\) を含む四つのバイベクトルが加わる。
この 10 次元空間が、PGA 版 DST の拡張ねじれ不整合の舞台となる（次章以降）。

表 \ref{tab:g311-grades} に \(G(3,1,1)\) の階数構造をまとめる。

\begin{table}[htbp]
  \centering
  \caption{\(G(3,1,1)\) の階数分解}
  \label{tab:g311-grades}
  \begin{tabular}{@{}clcl@{}}
    \toprule
    階数 & 次元 & 代表元 & 幾何的意味 \\
    \midrule
    0 & 1 & \(1\) & スカラー \\
    1 & 5 & \(e_\mu,\ e_4\) & ベクトル（時空＋射影方向） \\
    2 & 10 & \(e_\mu e_\nu,\ e_4 e_\mu\) & バイベクトル（平面・回転・並進生成子） \\
    3 & 10 & \(e_a e_b e_c\) & 三ベクトル（平面） \\
    4 & 5 & \(e_a e_b e_c e_d\) & 四ベクトル（点） \\
    5 & 1 & \(I_5\) & 5 次擬スカラー \\
    \bottomrule
  \end{tabular}
\end{table}

\subsection{擬スカラーの二重構造}

\(G(3,1,1)\) には、異なる役割を持つ二種類の擬スカラーが現れる。

\paragraph{4 次擬スカラー \(i\)}
\begin{equation}
  i = e_0 e_1 e_2 e_3
  \label{eq:g311-pseudoscalar-i}
\end{equation}
は \eqref{eq:cl31-pseudoscalar}・\eqref{eq:dst-pseudoscalar} と同じ元である。
\(i^2 = -1\) を満たし、DST の\textbf{双対写像} \(X \mapsto Xi\) の担い手となる——通常時空セクターと双対時空セクターを入れ替える。

\paragraph{5 次擬スカラー \(I_5\)}
\begin{equation}
  I_5 = e_0 e_1 e_2 e_3 e_4 = i\, e_4
  \label{eq:I5-pseudoscalar}
\end{equation}
は \(G(3,1,1)\) 全体の最高階成分である。
\(e_4^2 = 0\) より \(I_5^2 = (i e_4)^2 = i^2 e_4^2 = 0\) となり、\(I_5\) 自身も\textbf{ヌル}である——退化代数特有の性質である。
\(I_5\) は射影幾何における\textbf{代数的双対} \eqref{eq:algebraic-dual} の担い手となり、点と超平面の対応を与える（次節・次の節）。

\paragraph{\(e_4\) と \(i\) の可換性}
双対写像を \(G(3,1,1)\) へ拡張するうえで決定的なのは、\(e_4\) が \(i\) と可換であること \eqref{eq:e4-i-commute} である。
反可換関係 \(\{e_4, e_\mu\} = 0\) を 4 回用いれば
\begin{equation}
  e_4\, i
  = e_4\, e_0 e_1 e_2 e_3
  = (-1)^4\, e_0 e_1 e_2 e_3\, e_4
  = i\, e_4,
\end{equation}
すなわち \eqref{eq:e4-i-commute} が成り立つ。
したがって双対写像 \(X \mapsto Xi\) は \(\mathrm{Cl}(3,1)\) 部分代数上では第 II 部と同一に作用しつつ、\(e_4\) 方向は独立に保たれる——PGA 版 DST において「双対性を保ちながら並進を内在化する」代数的基盤がここにある。

\subsection{10 次元バイベクトル空間の概観}

\(G(3,1,1)\) の 2 次部分は 10 次元であり、次の三群に分類される（詳細は次章）。

\begin{itemize}[leftmargin=2em]
  \item \textbf{双曲生成子}（3 個、平方 \(+1\)）：\(iI = e_0 e_1,\ iJ = e_0 e_2,\ iK = e_0 e_3\)。
  第 II 部 \eqref{eq:hyperbolic-biv} と同じ。ブーストを担う。
  \item \textbf{巡回生成子}（3 個、平方 \(-1\)）：\(I = e_3 e_2,\ J = e_1 e_3,\ K = e_2 e_1\)。
  第 II 部 \eqref{eq:cyclic-biv} と同じ。空間回転を担う。
  \item \textbf{ヌル生成子}（4 個、平方 \(0\)）：\(N_\mu = e_4 \wedge e_\mu\)（\(\mu = 0,1,2,3\)）。
  \(e_4\) と各時空ベクトルの外積からなる。並進を担う——次章で本格導入する。
\end{itemize}

六つの双曲・巡回生成子は \(\mathfrak{so}(3,1) \oplus \mathfrak{so}(3,1)\) を張り、四つのヌル生成子は並進理想を与える。
全体として 10 次元バイベクトル空間はポアンカレ代数 \(\mathfrak{iso}(3,1)\) と一致する——これが PGA 版 DST の生成子構造の骨格である。

\section{点・直線・平面の多重ベクトル表現}

\subsection{射影幾何と PGA の基本対応}

射影幾何代数 \(G(3,1,1)\) では、幾何対象が多重ベクトル（とくにブレード）として符号化される。
本書および PGA 版論文は、\textbf{平面を 3-ブレード}として表す plane-based 規約を採用する——これは Lengyel--Gunn 系 PGA 文献の慣例に沿う。

第 2 章 \eqref{eq:k-blade} で定義した \(k\)-ブレードの考え方はそのまま \(G(3,1,1)\) に拡張される。
\(k\) 個の独立なベクトルの外積 \(a_1 \wedge a_2 \wedge \cdots \wedge a_k\) が \(k\)-ブレードを与える。
\(G(3,1,1)\) では \(k = 0,1,2,3,4,5\) まで現れるが、射影幾何の基本対象——点・直線・平面——は主に階数 1, 2, 3, 4 に対応する。

\subsection{平面：3-ブレード}

\textbf{平面}は 3-ブレードとして表される。
三つの独立なベクトルの外積が、有向平面（向きと大きさを持つ 2 次元部分空間）を与える。

\paragraph{通常時空側の平面}
\(e_4\) を含まない 3-ブレードは、通常時空生成子だけから構成される 2-平面を表す：
\begin{equation}
  P_{\mathrm{spatial}} = e_0 \wedge e_1 \wedge e_2.
  \label{eq:plane-spatial}
\end{equation}
これは \(e_0\)–\(e_1\)–\(e_2\) 方向に張られた有向 2-平面——\(tx\)-平面と \(ty\)-方向を含む時空 2-平面——に対応する。
一般に \(e_0 \wedge e_i \wedge e_j\)（\(i \neq j\)）は時空内の有向 2-平面を表す。
純空間的な例として \(e_1 \wedge e_2 \wedge e_3\) も 3-ブレードであるが、これは \(e_0\) 方向に垂直な空間 3-体積要素として \(\mathrm{Cl}(3,1)\) 内で解釈される——5 次元 PGA では 2-平面符号として読む。

\paragraph{ヌル因子を含む平面}
\(e_4\) を一因子として含む 3-ブレードは、射影幾何の\textbf{ヌル平面}を表す：
\begin{equation}
  P_{\mathrm{null}} = e_4 \wedge e_1 \wedge e_2.
  \label{eq:plane-null}
\end{equation}
\(e_4^2 = 0\) であるため、この 3-ブレードは退化計量下の「無限遠を含む平面」として機能する。
PGA 版論文では、平面 \(P\) と双対時空要素 \(D\) の関係——法線対応——が、このヌル平面構造を通じて記述される（第 3 節）。

\paragraph{平面の向き}
3-ブレード \(P = a \wedge b \wedge c\) は、三ベクトル \(a,b,c\) の順序によって向きが決まる。
\(a \wedge b \wedge c = -b \wedge a \wedge c\) のように、任意の二因子を入れ替えると符号が反転する——第 2 章の 2-ブレードと同様の性質である。

\subsection{直線：2-ブレード}

\textbf{直線}は 2-ブレードとして表される。
二つの独立なベクトルの外積が、有向直線（方向と位置を持つ 1 次元部分空間）を与える。

\paragraph{通常の直線}
\begin{equation}
  L = e_1 \wedge e_2
  \label{eq:line-example}
\end{equation}
は \(e_3\) 方向に垂直な \(xy\)-平面内の直線——厳密には、その平面内の有向 1-平面——を表す。
\(e_0 \wedge e_1\) は時間–\(x\) 方向の有向 2-平面、すなわち \(tx\)-平面内の「直線」に対応する。

\paragraph{ヌル直線}
\begin{equation}
  L_{\mathrm{null}} = e_4 \wedge e_\mu \qquad (\mu = 0,1,2,3)
  \label{eq:line-null}
\end{equation}
は \(e_4\) と時空ベクトル \(e_\mu\) からなる 2-ブレードである。
これは次章で\textbf{ヌル生成子} \(N_\mu = e_4 \wedge e_\mu\) として再登場する——並進方向そのものが、ヌル直線として幾何的に読める。
二平面の\textbf{交わり}（meet）として直線が現れることもある（次節）。

\subsection{点：1-ブレードと 4-ブレード}

射影幾何では、有限点と無限遠点を統一的に扱うために\textbf{斉次座標}が用いられる。
PGA では、この斉次化が代数内部に組み込まれている。

\paragraph{1-ブレードとしての点}
1-ブレード \(p = e_4 + x\, e_1 + y\, e_2 + z\, e_3\)（\(e_4\) 成分を含むベクトル）は、座標 \((x,y,z)\) を持つ有限点を表す——\(e_4\) 成分が「射影的な 1」に対応する。
純粋に \(e_4\) のみのベクトル \(e_4\) 自身は、無限遠方向——すべての有限直線が収束する方向——を表す。

\paragraph{4-ブレードとしての点（双対規約）}
5 次擬スカラー \(I_5\) \eqref{eq:I5-pseudoscalar} を用いた PGA 双対により、点は 4-ブレードとしても表せる：
\begin{equation}
  A_k^{\sim} := A_k\, I_5
  \qquad\text{（\(k\) 次ブレード \(A_k\) に対して）}.
  \label{eq:pga-dual}
\end{equation}
\(I_5^2 = 0\) であるため、通常の逆元 \(I_5^{-1}\) は存在しない——退化代数では \(I_5\) 自身が双対操作の担い手となる \eqref{eq:algebraic-dual} の退化版である。
この操作は \(k\) 次と \((5-k)\) 次を対応させる。
とくに 1-ブレード（ベクトル）と 4-ブレード（四重外積）が双対関係にあり、\textbf{点}と\textbf{超平面}が同一の代数構造の両面として現れる。

本書および PGA 版論文では、主に\textbf{平面ベース}（3-ブレード中心）で議論を進める。
点の 4-ブレード表現は join/meet の双対性を理解するうえで必要になるが、DST の物理的内容——ねじれ不整合、法線対応——は平面とその法線の関係として記述される。

\subsection{斉次座標とヌル構造}

\(e_4\) の導入は、射影幾何の斉次座標 \((w, x, y, z)\) を代数内部に内在化することに対応する。
有限点は \(w \neq 0\) の斉次座標 \((w, x, y, z) \sim (1, x/w, y/w, z/w)\) で表され、\(w = 0\) は無限遠点である。
代数では \(e_4\) 成分が \(w\) に対応し、\(e_4^2 = 0\) という退化性が「無限遠と有限点の区別」をヌル構造として符号化する。

並進がこの斉次座標にどう作用するか——すなわち \(N_\mu\) による有限シフト——は次章で扱う。
本章では、\(e_4\) が射影方向を担い、平面・直線・点がそれぞれ 3-, 2-, 1-（または 4-）ブレードとして符号化される、という幾何的対応を確立する。

\section{接合（join）と交わり（meet）の代数的実現}

\subsection{接合（join）}

射影幾何において、二つの部分空間の\textbf{接合}（join）とは、それらを同時に含む最小の部分空間である。
PGA では、独立なブレードの\textbf{外積}が join を与える。

\(k\)-ブレード \(A_k\) と \(\ell\)-ブレード \(B_\ell\) が独立（共通の非零ベクトルを含まない）なとき、
\begin{equation}
  \mathrm{join}(A_k, B_\ell) \;:=\; A_k \wedge B_\ell
  \label{eq:join-def}
\end{equation}
が \((k+\ell)\)-ブレードとして join を与える——ただし \(k+\ell \leq 5\) である必要がある。
以降、文脈が明らかなときは join を外積 \(A_k \wedge B_\ell\) と同一視してよい。

\paragraph{例：点と直線の join}
点 \(p = e_4 + e_1\)（\(x\) 軸上の原点）と直線 \(L = e_2 \wedge e_3\)（\(yz\)-平面内の直線）を考える。
\(p\) と \(L\) が独立なら
\[
  \mathrm{join}(p, L) = p \wedge L = (e_4 + e_1) \wedge e_2 \wedge e_3
  = e_4 \wedge e_2 \wedge e_3 + e_1 \wedge e_2 \wedge e_3
\]
は 3-ブレード——すなわち平面——を与える。

\paragraph{例：二直線の join}
同一平面内の二直線 \(L_1 = e_1 \wedge e_2\)、\(L_2 = e_1 \wedge e_3\) について
\[
  L_1 \wedge L_2 = e_1 \wedge e_2 \wedge e_3
\]
は \(e_1\) を法線とする平面を与える——二直線の join は、それらを含む平面である。

\subsection{交わり（meet）}

二つの部分空間の\textbf{交わり}（meet）とは、両方に共通する最大の部分空間である。
PGA では、meet は\textbf{双対}を介して定義される：
\begin{equation}
  A_k \vee B_\ell \;:=\; \bigl( A_k^{\sim} \wedge B_\ell^{\sim} \bigr)^{\sim},
  \label{eq:meet-def}
\end{equation}
ここで \(A_k^{\sim} = A_k\, I_5\) は \eqref{eq:pga-dual} の PGA 双対である。
外積 \(\wedge\) は join、記号 \(\vee\) は meet に用いる——文献によっては逆の規約もあるが、本書はこの取り方に統一する。

\paragraph{例：二平面の meet}
平面 \(P_1 = e_1 \wedge e_2 \wedge e_3\) と \(P_2 = e_4 \wedge e_1 \wedge e_2\) の meet を求める。
双対を計算すると meet は \(e_1 \wedge e_2\) 型の 2-ブレード——二平面の共通直線——になる。
具体的な \(I_5\) を用いた計算は演習に委ねる。

\paragraph{退化計量への注意}
\(e_4^2 = 0\) であるため、\(G(3,1,1)\) は\textbf{退化}（degenerate）代数である。
非退化代数 \(G(3)\) や \(G(3,1)\) で用いた標準的ホッジ双対 \(\star\) とは、符号や正規化が一致しない場合がある。
しかし射影幾何の join/meet 操作としては、\(I_5\) による双対 \eqref{eq:pga-dual} で十分に機能する——PGA 文献の標準的作法である。

\subsection{平面ベースのブラケット積}

PGA 版論文 Section 4 では、二つの 3-ブレード \(A, B\) に対する\textbf{平面ベースのブラケット積}が導入される。
幾何積の反対称部分と対称部分：
\begin{equation}
  A \mathbin{\underline{\vee}} B \;:=\; \tfrac{1}{2}(AB - BA),
  \qquad
  A \mathbin{\overline{\vee}} B \;:=\; \tfrac{1}{2}(AB + BA),
  \label{eq:bracket-def}
\end{equation}
を定義する。
ここで \(\mathbin{\underline{\vee}}\) と \(\mathbin{\overline{\vee}}\) は、標準 PGA 文献の anti-inner / anti-outer 積とは異なる——本書は論文の規約に従う \footnote{標準 PGA 文献（Lengyel, Dorst--Fontijne--Mann 等）の anti-inner / anti-outer 積は 5 次擬スカラーによる双対化で定義され、\eqref{eq:bracket-def} とは一般に異なる。}.

5 次元ベクトル空間では 6 次成分は存在しないため、3-ブレード同士の幾何積の階数内容は
\[
  A_3 B_3 = \langle A_3 B_3 \rangle_0 + \langle A_3 B_3 \rangle_2 + \langle A_3 B_3 \rangle_4
\]
に限定される。
直接計算により、各成分は \(A \leftrightarrow B\) の交換に対して次の対称性を持つ：
\[
  \langle A_3 B_3 \rangle_0 = \langle B_3 A_3 \rangle_0,
  \qquad
  \langle A_3 B_3 \rangle_2 = -\,\langle B_3 A_3 \rangle_2,
  \qquad
  \langle A_3 B_3 \rangle_4 = \langle B_3 A_3 \rangle_4.
\]
したがってブラケット積は階数を次のように分離する：
\begin{equation}
  A \mathbin{\underline{\vee}} B = \langle AB \rangle_2,
  \qquad
  A \mathbin{\overline{\vee}} B = \langle AB \rangle_0 + \langle AB \rangle_4.
  \label{eq:bracket-grades}
\end{equation}
交換子ブラケット \(\mathbin{\underline{\vee}}\) は 2 次部分——バイベクトル——を取り出し、二平面の\textbf{相対的な向き}（法線方向）を符号化する。
反交換子ブラケット \(\mathbin{\overline{\vee}}\) は 0 次と 4 次部分を取り出し、\textbf{射影の大きさ}と\textbf{並進的ずれ}を符号化する。

\subsection{双対時空と平面の法線}

PGA 版 DST の中心的結果の一つは、\textbf{双対時空}要素が 3-ブレードの法線として機能することである。
通常時空側の 3-ブレード \(P\) と双対時空セクターの要素 \(D\) に対し、交換子ブラケットは
\begin{equation}
  P \mathbin{\underline{\vee}} D = \kappa\,(e_4 \wedge n),
  \label{eq:normal-correspondence}
\end{equation}
の形をとる。
ここで \(n\) は \(P\) が表す平面の法線ベクトル（通常時空の意味で）、\(\kappa\) は \(P, D\) の正規化から決まる実係数である。
右辺 \(e_4 \wedge n\) はヌルバイベクトル \(N_\mu = e_4 \wedge e_\mu\) の線型結合——すなわち並進セクター——に属する。

この対応 \eqref{eq:normal-correspondence} は、双対時空が「平面の法線束」（normal bundle）として幾何的に読めることを示す。
ねじれ不整合の六自由度（双曲・巡回）と並進の四自由度（ヌル）の分解は、平面ベースのブラケット積によって幾何的に実現される——次章で構成する拡張バイベクトル \(\Omega_{\mathrm{biv}}^{(5)}\) の三分解と対応し、詳細な計算は第 IV 部に委ねる。

\paragraph{双対性との整合}
双対写像 \(X \mapsto Xi\) は \(\mathrm{Cl}(3,1)\) 部分代数上で通常時空と双対時空を入れ替え、ヌルセクター \(e_4 \wedge e_\mu\) には閉じた作用を与える \eqref{eq:e4-i-commute}。
したがって法線対応 \eqref{eq:normal-correspondence} は双対性の下でも保存される——10 次元生成子構造と \(G(3,1,1)\) の平面幾何が互換的である。

\section*{演習問題}

\begin{enumerate}[leftmargin=2em]
  \item 反可換関係 \(\{e_4, e_\mu\} = 0\) と \(e_4^2 = 0\) \eqref{eq:e4-definitions} から、任意の \(\mu\) について \(e_4 e_\mu e_4 = 0\) を示せ。

  \item \(i = e_0 e_1 e_2 e_3\) \eqref{eq:g311-pseudoscalar-i} について、\(\{e_4, e_\mu\} = 0\) を 4 回用いて \eqref{eq:e4-i-commute} を確認せよ。

  \item 3-ブレード \(P = e_4 \wedge e_1 \wedge e_2\) \eqref{eq:plane-null} と \(Q = e_1 \wedge e_2 \wedge e_3\) \eqref{eq:plane-spatial} について、それぞれの階数と \(e_4\) の有無による幾何的意味の違いを述べよ。

  \item 直線 \(L = e_1 \wedge e_2\) \eqref{eq:line-example} とベクトル \(e_3\) の join \(L \wedge e_3\) \eqref{eq:join-def} を計算せよ。
  結果が 3-ブレード \(e_1 \wedge e_2 \wedge e_3\) となることを確認せよ（空間 3-体積要素としての 3-ブレード）。

  \item 5 次擬スカラー \(I_5 = e_0 e_1 e_2 e_3 e_4\) \eqref{eq:I5-pseudoscalar} を用い、3-ブレード \(P = e_1 \wedge e_2 \wedge e_3\) の PGA 双対 \eqref{eq:pga-dual} \(P^{\sim} = P\, I_5\) を計算せよ。
  結果の階数と、点–超平面対応の直感を述べよ。

  \item 二つの 3-ブレード \(A, B\) の幾何積について、\(\langle AB \rangle_2 = -\,\langle BA \rangle_2\) であることを \eqref{eq:bracket-grades} の対称性から説明せよ。
  したがって \(A \mathbin{\underline{\vee}} B = \langle AB \rangle_2\) となる理由を述べよ。

  \item （発展）平面 \(P = e_4 \wedge e_1 \wedge e_2\) \eqref{eq:plane-null} と双対時空セクターの代表元 \(D = jI\)（\(j = e_0\)、\(I = e_3 e_2\)）について、交換子 \(PD - DP\) を計算せよ。
  \(\{e_4, e_\mu\} = 0\) と \(e_4^2 = 0\) を用い、結果が \(e_4 \wedge e_3\) 型のヌルバイベクトルに比例することを確認せよ。
  これは法線対応 \eqref{eq:normal-correspondence} の具体例である（論文 Appendix B 参照）。
\end{enumerate}

\chapter{並進とモーター}

前章では、ヌルベクトル \(e_4\) を加えた \(G(3,1,1)\) の射影構造——点・直線・平面のブレード表現、join/meet、平面ベースのブラケット積——を整えた。
本章では、その代数の\textbf{2 次部分}に焦点を当て、並進を担う\textbf{ヌル生成子} \(N_\mu\) の性質を確立し、\textbf{トランスレータ}と\textbf{拡張モーター} \(M = \exp(\Omega_{\mathrm{biv}}^{(5)})\) として並進を実装する。
第 3 章で予告した \eqref{eq:exp-null-truncation} の指数打ち切りと \eqref{eq:motor-sandwich} のサンドイッチ積を、PGA 版 DST の文脈で本格的に扱う。
平面ブラケットと法線対応 \eqref{eq:normal-correspondence} は前章で導入済みであり、本章では並進セクターがどのように有限シフトとして現れるかを中心に進める。
拡張不変量 \(J^{(5)}\) の詳細は次章に委ねる。

\section{\texorpdfstring{ヌル生成子 \(N_\mu = e_4 \wedge e_\mu\)}{ヌル生成子 N mu = e4 wedge e mu}}

\subsection{定義と明示表示}

前章 \eqref{eq:line-null} で直線として現れた 2-ブレードを、並進生成子として正式に定義する。
四つの\textbf{ヌル生成子}（null generator）は
\begin{equation}
  N_\mu := e_4 \wedge e_\mu = e_4 e_\mu
  \qquad (\mu = 0,1,2,3)
  \label{eq:null-generators}
\end{equation}
で与えられる。
ここで \(e_\mu\) は通常時空の四ベクトル \eqref{eq:cl31-basis} である。
外積記号 \(\wedge\) を省略した右辺 \(e_4 e_\mu\) は、\(e_4\) と \(e_\mu\) が反可換 \eqref{eq:e4-definitions} であるため、順序を入れ替えると符号が反転するだけで、2-ブレードとしての内容は同じである。

明示的に書けば
\begin{align}
  N_0 &= e_4 \wedge e_0 = e_4\, j, &
  N_1 &= e_4 \wedge e_1 = e_4\,(kI), \notag\\
  N_2 &= e_4 \wedge e_2 = e_4\,(kJ), &
  N_3 &= e_4 \wedge e_3 = e_4\,(kK)
  \label{eq:null-explicit}
\end{align}
となる。
各 \(N_\mu\) は \(e_4\) と時空方向 \(e_\mu\) が張る\textbf{ヌル直線} \eqref{eq:line-null} に対応し、並進の四方向を代数内部に符号化する。

\subsection{強消滅性質 \(N_\mu N_\nu = 0\)}

ヌル生成子の代数的特徴は、個々が平方ゼロであることだけではない。
\textbf{任意の組} \((\mu,\nu)\) に対して積がゼロになる——この\textbf{強消滅}（strong vanishing）が、並進指数の一次打ち切りを保証する。

\(\mu,\nu \in \{0,1,2,3\}\) に対し、反可換関係 \(\{e_4, e_\mu\} = 0\) と \(e_4^2 = 0\) \eqref{eq:e4-definitions} を用いると
\begin{equation}
  N_\mu N_\nu
  = (e_4 e_\mu)(e_4 e_\nu)
  = e_4\, e_\mu\, e_4\, e_\nu
  = -\,e_4\, e_4\, e_\mu e_\nu
  = -\,e_4^2\, e_\mu e_\nu
  = 0.
  \label{eq:null-strong-vanishing}
\end{equation}
\(\mu = \nu\) の場合は \(N_\mu^2 = 0\)（各生成子の平方ゼロ）、\(\mu \neq \nu\) の場合は異なる生成子同士の積もゼロ——いずれも \eqref{eq:null-strong-vanishing} に含まれる。
対角成分 \(N_\mu^2 = 0\) だけでは不十分であり、\textbf{すべての混合積} \(N_\mu N_\nu\) が消えることが、後述の指数打ち切り \eqref{eq:exp-truncation} の根拠となる。

\paragraph{計算例}
\(N_0 N_1 = (e_4 j)(e_4 kI)\) を計算する。
\(e_4\) は \(j\) とも \(kI\) とも反可換するから
\[
  N_0 N_1 = e_4 j e_4 kI = -e_4 e_4 j kI = -e_4^2 j kI = 0.
\]
四つの生成子は互いに\textbf{可換}でもある——\([N_\mu, N_\nu] = N_\mu N_\nu - N_\nu N_\mu = 0\)——したがって、ヌルセクターは\textbf{可換な理想}（abelian ideal）を形成する。

\subsection{ポアンカレ代数 \(\mathfrak{iso}(3,1)\) への埋め込み}

\(G(3,1,1)\) の 10 次元バイベクトル空間は、前章 \eqref{eq:g311-grade-dims} で概観した三群に分類される。

\begin{itemize}[leftmargin=2em]
  \item \textbf{双曲生成子}（3 個）：\(iI, iJ, iK\) \eqref{eq:hyperbolic-biv}。平方 \(+1\)。ブーストを担う。
  \item \textbf{巡回生成子}（3 個）：\(I, J, K\) \eqref{eq:cyclic-biv}。平方 \(-1\)。空間回転を担う。
  \item \textbf{ヌル生成子}（4 個）：\(N_\mu\) \eqref{eq:null-generators}。平方 \(0\)。並進を担う。
\end{itemize}

六つの双曲・巡回生成子の交換子は \(\mathfrak{so}(3,1)\) 内に閉じ、二つの \(\mathfrak{so}(3,1)\) コピーの直和 \(\mathfrak{so}(3,1) \oplus \mathfrak{so}(3,1)\) を張る——第 II 部 \eqref{eq:hyperbolic-biv}–\eqref{eq:cyclic-biv} の構造の拡張である。
四つのヌル生成子は可換理想をなし、ポアンカレ代数の並進生成子 \(P_\mu\) に対応する。
全体として
\begin{equation}
  \dim \mathfrak{iso}(3,1)
  = \underbrace{3 + 3}_{\mathfrak{so}(3,1) \oplus \mathfrak{so}(3,1)}
  + \underbrace{4}_{\text{並進理想}}
  = 10
  \label{eq:poincare-iso}
\end{equation}
次元のリー代数 \(\mathfrak{iso}(3,1)\) が、\(G(3,1,1)\) の 2 次部分に自然に実現される。
並進は座標変換として外付けするのではなく、代数のヌルセクターから生じる——PGA 版 DST が \(G(3,1,1)\) を選ぶ中心的理由である。

\subsection{なぜバイベクトル（2 次）なのか}

並進生成子を 1 次のヌルベクトルではなく 2 次のヌルバイベクトル \(N_\mu = e_4 \wedge e_\mu\) として実装する理由は、\textbf{反転}（reversion）の性質にある。

バイベクトルは反転奇性 \eqref{eq:biv-rev-odd} を持つ：
\begin{equation}
  \widetilde{N_\mu} = -N_\mu.
  \label{eq:null-rev-odd}
\end{equation}
1 次のベクトル \(v\) は \(\widetilde{v} = v\)（反転偶）であるから、もし \(e_4\) 自身を並進生成子に使おうとすると、後述のユニタリ性 \(M\widetilde{M} = 1\) が成り立たなくなる \footnote{Selig 以来の PGA 文献でも、ユークリッド・ミンコフスキー双方で並進はヌル\textbf{バイベクトル}として符号化される——第 2 章 \eqref{eq:biv-rev-odd} の脚注参照。}.
\(N_\mu\) が 2 次であること——\(e_4\) と \(e_\mu\) の\textbf{外積}——が、モーターの因果的伝播を保証する代数的根拠である。

\subsection{双対写像下でのヌルセクターの閉性}

双対写像 \(X \mapsto Xi\) は \(\mathrm{Cl}(3,1)\) 部分代数上で通常時空と双対時空を入れ替える。
\(e_4\) は擬スカラー \(i = e_0 e_1 e_2 e_3\) と可換 \eqref{eq:e4-i-commute} であるから、ヌル生成子は双対の下でもヌルセクター内に留まる：
\begin{equation}
  N_\mu\, i = e_4 e_\mu i = e_4 (e_\mu i),
\end{equation}
右辺は再び \(e_4 \wedge (\text{ベクトル})\) の形——すなわち \(e_4 \wedge (e_\mu i)\)——であり、4 次元ヌル部分空間に属する。
したがって 10 次元生成子空間の三分解（双曲・巡回・ヌル）は双対性の下でも保存される。
双対写像とヌルセクターの詳細な相互作用は次章「双対性とトーション関連の構造」で扱う。

\section{並進をモーターとして実現する方法}

\subsection{トランスレータ}

並進パラメータ \(a^\mu\)（\(\mu = 0,1,2,3\)）に対し、\textbf{トランスレータ}（translator）を
\begin{equation}
  T(a) := \tfrac{1}{2}\, a^\mu N_\mu
  \label{eq:translator}
\end{equation}
と定義する（アインシュタインの和の規約を用いる）。
\(T(a)\) はヌルセクターだけから構成される 2 次多重ベクトルである。

強消滅 \eqref{eq:null-strong-vanishing} より \(T(a)^2 = 0\) である：
\begin{equation}
  T(a)^2
  = \tfrac{1}{4}\, a^\mu a^\nu\, N_\mu N_\nu
  = 0.
  \label{eq:translator-square-zero}
\end{equation}
したがって指数写像 \eqref{eq:exp-biv} は 1 次で打ち切られ、
\begin{equation}
  \exp\!\bigl(T(a)\bigr)
  = 1 + T(a)
  = 1 + \tfrac{1}{2}\, a^\mu N_\mu.
  \label{eq:exp-truncation}
\end{equation}
並進は高次補正なしに\textbf{有限のシフト}として実現される——曲率や追加の高次項を導入せず、純粋な平行移動として現れる。
第 3 章 \eqref{eq:exp-null-truncation} で予告した公式の、パラメータ付きの一般形が \eqref{eq:exp-truncation} である。

\(\eta_{\mu\nu} = \mathrm{diag}(-1,+1,+1,+1)\) をミンコフスキー計量とすると、\textbf{光的並進}は
\begin{equation}
  \eta_{\mu\nu}\, a^\mu a^\nu = 0
  \label{eq:light-cone-translation}
\end{equation}
を満たすパラメータ \(a^\mu\) に対応する——第 IV 部で光子の移動として再登場する。

\subsection{拡張バイベクトル \(\Omega_{\mathrm{biv}}^{(5)}\) の三分解}

PGA 版 DST の\textbf{拡張ねじれ不整合}は、10 次元バイベクトル空間上の元
\begin{equation}
  \Omega_{\mathrm{biv}}^{(5)}
  = \underbrace{\sum_{a=1}^{3} \Bigl( \tfrac{\alpha_a}{2}\, B_a^{+} + \tfrac{\beta_a}{2}\, B_a^{-} \Bigr)}_{\Omega_{\mathrm{torsion}}\ :\ \text{ねじれ（6D）}}
  \;+\;
  \underbrace{\sum_{\mu=0}^{3} \tfrac{\lambda^\mu}{2}\, N_\mu}_{\Omega_{\mathrm{trans}}\ :\ \text{並進（4D）}}
  \label{eq:Omega-decomp}
\end{equation}
として書ける。
ここで \(B_a^{+}\) は双曲生成子 \((iI, iJ, iK)\)、\(B_a^{-}\) は巡回生成子 \((I, J, K)\)、\(N_\mu\) はヌル生成子 \eqref{eq:null-generators} である。

係数 \(\alpha_a, \beta_a\) は通常時空と双対時空の間の\textbf{ねじれ不整合}をパラメータ化し、四係数 \(\lambda^\mu\) は\textbf{並進的不整合}をローレンツ 4 ベクトルとして符号化する。
光的並進は \(\eta_{\mu\nu}\lambda^\mu\lambda^\nu = 0\) \eqref{eq:light-cone-translation} で特徴づけられる。

\subsection{モーターの構成と分解 \(M = R\,T\)}

拡張モーター（extended motor）は
\begin{equation}
  M = \exp\!\bigl(\Omega_{\mathrm{biv}}^{(5)}\bigr)
  \label{eq:extended-motor}
\end{equation}
で定義される。
第 3 章のローター \eqref{eq:rotor-def} と同型の構造を持ち、ねじれ・並進の両方を単一の指数元から生成する。

ねじれ部分と並進部分は一般に\textbf{非可換}である——双曲・巡回生成子とヌル生成子の交換子 \([B_a^{\pm}, N_\mu]\) は一般に非零——が、ポアンカレ代数の半直積構造 \eqref{eq:poincare-iso} により、モーターは標準的に
\begin{equation}
  M = R\, T,
  \qquad
  R := \exp(\Omega_{\mathrm{torsion}}),
  \qquad
  T := \exp(\Omega'_{\mathrm{trans}}),
  \label{eq:motor-factor}
\end{equation}
と分解できる。
ここで \(\Omega'_{\mathrm{trans}} = \sum_\mu (\lambda'^{\mu}/2)\, N_\mu\) は、\(R\) による共役によって再パラメータ化された並進生成子であり、\(\lambda'^{\mu}\) は元の \(\lambda^\mu\) から \(R\) の作用で得られる。
トランスレータ \(T\) 自体は \eqref{eq:exp-truncation} どおりヌルバイベクトルの\textbf{一次多項式}のままである——並進的シフトに高次補正は入らない。

\subsection{ユニタリ性 \(M\widetilde{M} = 1\)}

PGA 版 DST において、幾何変換が因果的であるための条件は、モーターの\textbf{ユニタリ性}
\begin{equation}
  M\,\widetilde{M} = 1
  \label{eq:motor-unitarity}
\end{equation}
である——第 3 章 \eqref{eq:rotor-def} の \(R\widetilde{R} = 1\) の拡張版である。

\paragraph{ねじれローター側}
すべてのバイベクトル生成子（双曲・巡回・ヌル）は反転奇性 \eqref{eq:biv-rev-odd}・\eqref{eq:null-rev-odd} を満たすから、
\begin{equation}
  R\,\widetilde{R}
  = \exp(\Omega_{\mathrm{torsion}})\,\exp(\widetilde{\Omega_{\mathrm{torsion}}})
  = \exp(\Omega_{\mathrm{torsion}})\,\exp(-\Omega_{\mathrm{torsion}})
  = 1.
\end{equation}

\paragraph{トランスレータ側}
\eqref{eq:exp-truncation} と \eqref{eq:null-rev-odd} を用いると
\begin{align}
  T\,\widetilde{T}
  &= \Bigl(1 + \sum_\mu \tfrac{\lambda^\mu}{2}\, N_\mu\Bigr)
     \Bigl(1 + \sum_\nu \tfrac{\lambda^\nu}{2}\,\widetilde{N_\nu}\Bigr) \notag\\
  &= \Bigl(1 + \sum_\mu \tfrac{\lambda^\mu}{2}\, N_\mu\Bigr)
     \Bigl(1 - \sum_\nu \tfrac{\lambda^\nu}{2}\, N_\nu\Bigr) \notag\\
  &= 1 - \underbrace{\Bigl(\sum_\mu \tfrac{\lambda^\mu}{2}\, N_\mu\Bigr)^{\!2}}_{=\,0\ \text{by \eqref{eq:null-strong-vanishing}}}
  = 1.
  \label{eq:null-unitarity}
\end{align}
強消滅 \eqref{eq:null-strong-vanishing} がここで決定的な役割を果たす——対角の \(N_\mu^2 = 0\) だけでは \eqref{eq:null-unitarity} の展開中間項を消せない。

\paragraph{全体}
反転は反準同型 \(\widetilde{RT} = \widetilde{T}\,\widetilde{R}\) であるから、\eqref{eq:motor-factor} より
\begin{equation}
  M\,\widetilde{M}
  = R\,T\,\widetilde{T}\,\widetilde{R}
  = R\,\widetilde{R}
  = 1.
\end{equation}
これが拡張モーターのユニタリ性 \eqref{eq:motor-unitarity} である。

\subsection{サンドイッチ積と計量の保存}

ユニタリ性 \eqref{eq:motor-unitarity} により、モーターによる変換は第 3 章 \eqref{eq:motor-sandwich} と同型の\textbf{サンドイッチ積}
\begin{equation}
  X' = M\,X\,\widetilde{M}
  \label{eq:motor-sandwich-chapter}
\end{equation}
で与えられる。
\(G(3,1,1)\) の退化計量——とくに通常時空ベクトル \eqref{eq:spacetime-vector} のノルム \eqref{eq:spacetime-norm}——はこの作用の下で保存される。
光的成分を含むベクトルのノルムも不変であり、並進パラメータが光円錐 \eqref{eq:light-cone-translation} 上にあるとき、変換は\textbf{光的伝播}として記述される。

前章の法線対応 \eqref{eq:normal-correspondence} と合わせると、双対時空要素が平面の法線として機能し、並進セクター \(N_\mu\) が平面に沿った有限シフトを与える——ねじれ不整合（六自由度）と並進（四自由度）が、単一のモーター \(M\) のサンドイッチ積 \eqref{eq:motor-sandwich-chapter} の枠内に統一される。

\section{剛体運動の統一的記述}

\subsection{モーターとポアンカレ群}

剛体運動学において、回転と並進は従来別々に扱われてきた。
射影幾何代数では、両者を\textbf{モーター}（motor）\(M = \exp(\Omega)\) として統一する——第 3 章の「ベルソル → ローター → 指数写像 → サンドイッチ積」の鎖 \eqref{eq:versor-action}–\eqref{eq:motor-sandwich} の \(G(3,1,1)\) 版である。

ユークリッド 3 次元空間の剛体運動群 \(\mathrm{SE}(3)\) では、6 次元バイベクトル（3 回転 + 3 並進）からなるモーターが回転と平行移動を同時に記述する。
\(G(3,1,1)\) ではこの構造が\textbf{ミンコフスキー版}に拡張され、10 次元バイベクトル空間 \eqref{eq:poincare-iso} がポアンカレ群 \(\mathrm{ISO}(3,1)\)——ローレンツ変換と時空並進の半直積——のリー代数 \(\mathfrak{iso}(3,1)\) を実現する。

\begin{table}[htbp]
  \centering
  \caption{剛体運動と PGA 版 DST の対応}
  \label{tab:motor-comparison}
  \begin{tabular}{@{}lll@{}}
    \toprule
    & ユークリッド \(\mathrm{SE}(3)\) & ミンコフスキー \(\mathrm{ISO}(3,1)\) \\
    \midrule
    代数 & \(G(3)\) 型 PGA & \(G(3,1,1)\) \\
    バイベクトル次元 & 6 & 10 \\
    回転・ブースト & 3 回転 & 3 回転 + 3 ブースト \\
    並進 & 3 並進 & 4 並進 \\
    モーター & \(M = \exp(\Omega_{6})\) & \(M = \exp(\Omega_{\mathrm{biv}}^{(5)})\) \\
    ユニタリ性 & \(M\widetilde{M} = 1\) & \eqref{eq:motor-unitarity} \\
    \bottomrule
  \end{tabular}
\end{table}

\subsection{ねじれと並進の単一オブジェクト}

PGA 版 DST において、拡張バイベクトル \eqref{eq:Omega-decomp} は、通常時空と双対時空の間の\textbf{ねじれ不整合}（六自由度）と、\textbf{並進的不整合}（四自由度）を単一の 10 次元オブジェクトに統合する。
六つの双曲・巡回生成子は親理論の torsional rotor を担い、四つのヌル生成子は並進セクターを担う——いずれも同じ指数写像 \eqref{eq:extended-motor} から生じ、同じサンドイッチ積 \eqref{eq:motor-sandwich-chapter} で伝播する。

光の自由度とねじれ不整合は、別個の概念としてではなく、10 次元生成子の異なる成分として等価に扱われる。
並進パラメータ \(\lambda^\mu\) が光円錐 \eqref{eq:light-cone-translation} 上にあるとき、並進は\textbf{零}（null）であり、ミンコフスキー計量の下で「長さゼロの移動」として特徴づけられる——これが PGA 版 DST における光子の光速移動の代数的基盤である。

\subsection{曲率なしの有限シフト}

トランスレータ \eqref{eq:exp-truncation} の一次打ち切りは、並進が\textbf{曲率テンソルを介さず}有限のシフトとして現れることを意味する。
背景連続体上の座標パッチを仮定するのではなく、粒子内在のモーター \(M\) が、ねじれと並進を同時に実行する——これが DST の代数的立場である。
Schwarzschild 解など具体的な物理帰結の回復は第 V 部で論じる。

\subsection{次章への接続：拡張不変量 \(J^{(5)}\)}

10 次元生成子空間上の二次形式として、PGA 版 DST は拡張不変量 \(J^{(5)}\) を導入する。
ポアンカレ代数のキリング形式は並進生成子に対して\textbf{退化}する——\(B_{\mathrm{Killing}}(P_\mu, P_\nu) = 0\)——ため、並進セクターにはミンコフスキー計量 \(\eta_{\mu\nu}\) を独立に付加する必要がある：
\begin{equation}
  J^{(5)}
  = \tfrac{1}{16}\, B_{\mathrm{Killing}}(\Omega_{\mathrm{torsion}}, \Omega_{\mathrm{torsion}})
  + \tfrac{1}{2}\, \eta_{\mu\nu}\, \lambda^\mu \lambda^\nu.
  \label{eq:J5-preview}
\end{equation}
第 2 項は光円錐 \eqref{eq:light-cone-translation} 上で正確に消え、光的並進が零の二次形式を持つことが組み込まれている。
\(J^{(5)}\) の導出、キリング形式の役割、ねじれ不整合との結合は次章「双対性とトーション関連の構造」で体系的に扱う。

\section*{演習問題}

\begin{enumerate}[leftmargin=2em]
  \item 反可換関係 \(\{e_4, e_\mu\} = 0\) と \(e_4^2 = 0\) \eqref{eq:e4-definitions} を用い、\(N_0 N_1 = 0\) \eqref{eq:null-strong-vanishing} を明示的に計算せよ。
  各ステップで \(e_4\) を左に移動する際の符号変化に注意すること。

  \item トランスレータ \(T(a) = \tfrac{1}{2} a^\mu N_\mu\) \eqref{eq:translator} について、\(T(a)^2 = 0\) \eqref{eq:translator-square-zero} を示し、そこから \eqref{eq:exp-truncation} が従うことを確認せよ。

  \item ヌル生成子 \(N_\mu = e_4 e_\mu\) \eqref{eq:null-generators} について、反転 \eqref{eq:biv-rev-odd} から \(\widetilde{N_\mu} = -N_\mu\) \eqref{eq:null-rev-odd} を導け。

  \item トランスレータ \(T = 1 + \sum_\mu (\lambda^\mu/2)\, N_\mu\) \eqref{eq:exp-truncation} について、\(\widetilde{N_\mu} = -N_\mu\) と \eqref{eq:null-strong-vanishing} を用いて \(T\widetilde{T} = 1\) \eqref{eq:null-unitarity} を直接計算せよ。

  \item モーターが \(M = RT\) \eqref{eq:motor-factor} と分解され、\(R\widetilde{R} = T\widetilde{T} = 1\) かつ反転が反準同型 \(\widetilde{RT} = \widetilde{T}\,\widetilde{R}\) であるとき、\(M\widetilde{M} = 1\) \eqref{eq:motor-unitarity} が従うことを示せ。

  \item 拡張バイベクトル \eqref{eq:Omega-decomp} の並進係数 \(\lambda^\mu = (c,\, 1,\, 0,\, 0)\)（\(c\) は光速）について、\(\eta_{\mu\nu}\lambda^\mu\lambda^\nu = 0\) \eqref{eq:light-cone-translation} となることを確認せよ。
  このパラメータが「光的並進」に対応することを、ミンコフスキー計量 \((-,+,+,+)\) の文脈で説明せよ。

  \item \(e_4\) と \(i = e_0 e_1 e_2 e_3\) \eqref{eq:g311-pseudoscalar-i} の可換性 \eqref{eq:e4-i-commute} を用い、\(N_0 i = e_4 (j i)\) が再び \(e_4 \wedge (\text{ベクトル})\) の形になることを示せ。
  したがって双対写像 \(X \mapsto Xi\) はヌルセクターを閉じたままにすることを述べよ。

  \item （発展）拡張不変量 \eqref{eq:J5-preview} の第 2 項 \(\tfrac{1}{2}\eta_{\mu\nu}\lambda^\mu\lambda^\nu\) が、\eqref{eq:light-cone-translation} のときゼロになることを確認せよ。
  この事実が「光子の移動は並進二次形式が零である」ことを意味することを、次章の議論の予習として述べよ。
\end{enumerate}

\chapter{双対性とトーション関連の構造}

前章では、ヌル生成子 \(N_\mu\) から拡張モーター \(M = \exp(\Omega_{\mathrm{biv}}^{(5)})\) を構成し、並進とねじれを単一のユニタリ変換として統合した。
本章では、PGA 版 DST の代数的骨格——\textbf{双対写像}、\textbf{極性}、\textbf{キリング形式}——を体系的に整え、拡張不変量 \(J^{(5)}\) を前章 \eqref{eq:J5-preview} の予告から正式定義へ昇格させる。
さらに、ねじれ不整合（torsional mismatch）と並進的不整合が、10 次元バイベクトル空間のどの部分に「住むか」を幾何的に位置づけ、第 IV 部の二重時空理論への橋を架ける。
モーターの構成や \(N_\mu\) の基本性質は前章で済んでいるため、本章ではそれらを前提として双対性と二次不変量の側面に焦点を当てる。

\section{双対写像と極性}

\subsection{双対写像 \(X \mapsto Xi\) の再確認}

DST の親理論および PGA 版論文が前提とする\textbf{双対写像}（duality map）は、4 次擬スカラー
\begin{equation}
  i = e_0 e_1 e_2 e_3
  \label{eq:duality-pseudoscalar}
\end{equation}
による右乗算
\begin{equation}
  X \;\longmapsto\; X\, i
  \label{eq:duality-map}
\end{equation}
として定義される \eqref{eq:cl31-pseudoscalar}。
この写像は \(\mathrm{Cl}(3,1)\) 部分代数上で、通常時空セクター（双曲生成子 \(\{iI,iJ,iK\}\) が担うブースト中心）と双対時空セクター（巡回生成子 \(\{I,J,K\}\) が担う回転中心）を入れ替える——第 II 部 \eqref{eq:hyperbolic-biv}–\eqref{eq:cyclic-biv} の物理的分割を代数内部で実現する。

\(G(3,1,1)\) へ拡張するとき、決定的なのは \(e_4\) が \(i\) と\textbf{可換}であること \eqref{eq:e4-i-commute} である。
反可換関係 \(\{e_4, e_\mu\} = 0\) を 4 回用いれば
\begin{equation}
  e_4\, i = e_4\, e_0 e_1 e_2 e_3 = (-1)^4\, e_0 e_1 e_2 e_3\, e_4 = i\, e_4,
\end{equation}
すなわち \eqref{eq:e4-i-commute} が成り立つ。
したがって双対写像 \eqref{eq:duality-map} は \(\mathrm{Cl}(3,1)\) 上では親理論と同一に作用しつつ、射影方向 \(e_4\) は独立なヌル方向として不変のままである——PGA 版 DST において「双対性を保ちながら並進を内在化する」という主張の代数的根拠がここにある。

\subsection{10 次元生成子空間への作用}

10 次元バイベクトル空間 \eqref{eq:poincare-iso} の三分解（双曲・巡回・ヌル）が双対写像 \eqref{eq:duality-map} の下でどう変形するかを、生成子ごとに見る。

\paragraph{双曲セクターと巡回セクターの入れ替え}
双曲生成子 \(B_a^{+} \in \{iI, iJ, iK\}\) と巡回生成子 \(B_a^{-} \in \{I, J, K\}\) は、双対写像の下で互いに対応する。
例えば \(iI = e_0 e_1\) に対し、
\begin{equation}
  (iI)\, i = e_0 e_1 e_0 e_1 e_2 e_3
  = -e_0^2 e_1^2 e_2 e_3
  = -(-1)(+1)\, e_2 e_3
  = -\, e_2 e_1
  = K,
  \label{eq:duality-hyperbolic-example}
\end{equation}
ここで \(K = e_2 e_1\) \eqref{eq:cyclic-biv} である。
同様に \(iJ \mapsto J\)、\(iK \mapsto I\)（符号は基底の取り方に依存するが、双曲セクター全体と巡回セクター全体が一対一対応する）。
これが「通常時空と双対時空が入れ替わる」という主張の生成子レベルでの内容である——前章 \eqref{eq:Omega-decomp} のねじれ部分 \(\Omega_{\mathrm{torsion}}\) は、双曲成分 \(\alpha_a\) と巡回成分 \(\beta_a\) の両方を含むが、双対写像はこの二成分を相互に写す。

\paragraph{ヌルセクターの閉性}
ヌル生成子 \(N_\mu = e_4 \wedge e_\mu\) \eqref{eq:null-generators} に対し、
\begin{equation}
  N_\mu\, i = e_4 e_\mu i = e_4 (e_\mu i),
  \label{eq:null-duality-closure}
\end{equation}
右辺は \(e_\mu i\) が \(\mathrm{Cl}(3,1)\) 内のベクトル（双対セクターの元）であるから、再び \(e_4 \wedge (\text{ベクトル})\) の形——すなわち \(e_4 \wedge (e_\mu i)\)——をとる。
これは 4 次元ヌル部分空間 \(\mathrm{span}\{N_0,N_1,N_2,N_3\}\) 内に留まる \eqref{eq:null-explicit}。
したがって双対写像はヌルセクターを\textbf{自身へ}写し、10 次元生成子空間の三分解（3 + 3 + 4）は双対性の下でも保存される \eqref{eq:poincare-iso}。

前章ではこの閉性を事実として述べたが、本章では生成子レベルで計算を確認した。
双対写像とヌルセクターの相互作用——とくに法線対応 \eqref{eq:normal-correspondence} との整合——は次節の極性と合わせて理解する。

\subsection{極性の二層：PGA 双対と DST 双対}

\(G(3,1,1)\) には、役割の異なる二種類の「双対」が共存する。

\paragraph{PGA 極性（射影双対）}
5 次擬スカラー \(I_5 = e_0 e_1 e_2 e_3 e_4 = i\, e_4\) \eqref{eq:I5-pseudoscalar} による PGA 双対 \eqref{eq:pga-dual} は、射影幾何における\textbf{極性}（polarity）——点と超平面、\(k\) 次ブレードと \((5-k)\) 次ブレードの対応——を与える。
\(I_5^2 = 0\) であるため逆元 \(I_5^{-1}\) は存在せず、\(I_5\) 自身が双対操作の担い手となる——退化代数 \(G(3,1,1)\) 特有の性質である \eqref{eq:algebraic-dual}。
join/meet \eqref{eq:join-def}–\eqref{eq:meet-def} や点–超平面対応は、PGA 双対 \eqref{eq:pga-dual} に基づく。

\paragraph{DST 双対（セクター入れ替え）}
一方、\eqref{eq:duality-map} の \(X \mapsto Xi\) は、\(\mathrm{Cl}(3,1)\) 部分代数上で通常時空と双対時空の\textbf{セクター}を入れ替える——物理的内容（ブースト中心 vs 回転中心）に対応する双対である \eqref{eq:dst-pseudoscalar}。

\paragraph{二層の両立}
両者は独立ではない。
\(I_5 = i\, e_4\) かつ \(e_4 i = i e_4\) \eqref{eq:e4-i-commute} であるから、射影極性と DST 双対は同じ代数の異なる面として共存する。
とくに \(e_4\) が \(i\) と可換であること——\(e_4\) が DST 双対の「通り道」に乗らないこと——が、並進方向を保ちつつ通常／双対セクターの入れ替えを可能にする。
PGA 版 DST では、平面ベースの議論 \eqref{eq:normal-correspondence} が DST 双対（法線＝双対時空）を担い、join/meet が PGA 極性を担う——二層が幾何と物理を同時に記述する。

\subsection{法線対応の双対共変性}

前章で導入した法線対応 \eqref{eq:normal-correspondence} は、双対写像 \eqref{eq:duality-map} の下でも形を保つ——PGA 版論文 Section 4 の\textbf{双対共変性}（duality covariance）である。

通常時空側の 3-ブレード \(P\) と双対時空セクターの元 \(D\) に対し、
\begin{equation}
  P \mathbin{\underline{\vee}} D = \kappa\,(e_4 \wedge n),
\end{equation}
ここで \(n\) は \(P\) の法線ベクトル、\(\kappa\) は正規化定数である \eqref{eq:normal-correspondence}。
双対写像を両辺に適用すると、\(P \mapsto P i\)、\(D \mapsto D i\)、\(n \mapsto n i\)（\(n\) は通常時空ベクトル）であり、
\begin{equation}
  (P i) \mathbin{\underline{\vee}} (D i)
  = \kappa'\, \bigl(e_4 \wedge (n i)\bigr),
  \label{eq:normal-duality-covariance}
\end{equation}
右辺の \(e_4 \wedge (n i)\) は再びヌルバイベクトル——双対写像後の法線方向を担う \eqref{eq:null-duality-closure}。
したがって法線対応は双対セクター間でも同型に成立し、10 次元生成子構造と \(G(3,1,1)\) の平面幾何が互換的である \eqref{eq:bracket-grades}。

\paragraph{計算例}
平面 \(P = e_4 \wedge e_1 \wedge e_2\) \eqref{eq:plane-null} と双対時空セクターの代表元 \(D = jI\)（\(j = e_0\)、\(I = e_3 e_2\)）について、交換子ブラケットを計算する。
\(\{e_4, e_\mu\} = 0\) と \(e_4^2 = 0\) \eqref{eq:e4-definitions} を用いると、
\[
  P \mathbin{\underline{\vee}} D
  = \tfrac{1}{2}(PD - DP)
  \propto e_4 \wedge e_3,
\]
すなわち \(N_3 = e_4 \wedge e_3\) 型のヌルバイベクトルに比例する——法線対応 \eqref{eq:normal-correspondence} の具体例である（論文 Appendix B 参照）。
双対写像 \(X \mapsto Xi\) を適用すると、\(P\) は双対セクターの 3-ブレードへ、\(D\) は通常セクターへ写され、\eqref{eq:normal-duality-covariance} の形が保たれる。

\section{キリング形式の役割}

\subsection{リー代数上の内積としてのキリング形式}

半単純リー代数 \(\mathfrak{g}\) 上の\textbf{カルタン--キリング形式}（Cartan--Killing form）は、
\begin{equation}
  B_{\mathrm{Killing}}(X, Y)
  := \mathrm{Tr}(\mathrm{ad}_X \circ \mathrm{ad}_Y),
  \label{eq:killing-def}
\end{equation}
で定義される対称双線形形式である。
ここで \(\mathrm{ad}_X(Y) = [X,Y]\) は随伴表現、\(\mathrm{Tr}\) は適切な基底での行列のトレースである。
キリング形式はリー代数の構造——とくに半単純性——を二次形式で測る道具であり、DST 親理論ではねじれ不整合スカラー \(J\) の定義に現れる \footnote{親理論では \(B_{\mathrm{Killing}}(X,Y) = 4\,\mathrm{Tr}(XY)\) 型の正規化が用いられることが多い。PGA 版論文 \eqref{eq:J5-invariant} の係数 \(1/16\) はこの正規化に合わせている。}.

\(G(3,1,1)\) の 10 次元バイベクトル空間 \eqref{eq:poincare-iso} では、六つの双曲・巡回生成子が \(\mathfrak{so}(3,1) \oplus \mathfrak{so}(3,1)\) を張る。
この半単純部分に対するキリング形式は、ねじれ部分 \(\Omega_{\mathrm{torsion}}\) \eqref{eq:Omega-decomp} の二次形式
\begin{equation}
  B_{\mathrm{Killing}}(\Omega_{\mathrm{torsion}}, \Omega_{\mathrm{torsion}})
  \label{eq:killing-torsion}
\end{equation}
として機能し、親理論のねじれスカラー \(J\) を再現する——第 IV 部で \(R_{\mathrm{usual}}\) と \(R_{\mathrm{dual}}\) の相対ずれとして物理的内容を与える。

\subsection{ポアンカレ代数での退化}

しかし、ポアンカレ代数 \(\mathfrak{iso}(3,1)\) 全体にキリング形式 \eqref{eq:killing-def} を適用すると、\textbf{退化}（degenerate）する。
並進生成子 \(P_\mu\)（代数上はヌルバイベクトル \(N_\mu\) に対応 \eqref{eq:null-generators}）に対し
\begin{equation}
  B_{\mathrm{Killing}}(P_\mu, P_\nu) = 0
  \label{eq:killing-degeneracy}
\end{equation}
が成り立つ——これはリー代数論の標準的事実である \footnote{可解理想（並進部分）が存在する代数では、キリング形式はその理想に対して零になる。}.
したがってキリング形式\textbf{単独}では、並進的不整合 \(\lambda^\mu\) を検出する二次不変量を与えられない。
前章 \eqref{eq:J5-preview} で予告した拡張が必要な理由がここにある。

\subsection{拡張不変量 \(J^{(5)}\) の正式定義}

PGA 版 DST では、半単純セクター上のキリング形式に、並進セクター上のミンコフスキー計量 \(\eta_{\mu\nu} = \mathrm{diag}(-1,+1,+1,+1)\) を独立に付加し、10 次元生成子空間上の\textbf{非退化}な二次不変量を構成する。
前章 \eqref{eq:J5-preview} の予告を、ここで正式に定義する：
\begin{equation}
  \boxed{\;
  J^{(5)}
  := \tfrac{1}{16}\, B_{\mathrm{Killing}}(\Omega_{\mathrm{torsion}}, \Omega_{\mathrm{torsion}})
  + \tfrac{1}{2}\, \eta_{\mu\nu}\, \lambda^\mu \lambda^\nu,
  \;}
  \label{eq:J5-invariant}
\end{equation}
ここで \(\Omega_{\mathrm{torsion}}\) と \(\lambda^\mu\) は拡張バイベクトル \eqref{eq:Omega-decomp} のねじれ係数 \((\alpha_a, \beta_a)\) と並進係数である。

第 1 項は親理論のねじれスカラー \(J\) を再現する——通常時空と双対時空の相対的ずれの大きさを測る。
第 2 項は並進的不整合 \(\lambda^\mu\) のローレンツ不変な二次測度を与える。
二つの寄与は\textbf{結合しない}——\(\mathfrak{iso}(3,1)\) 上のキリング形式は半単純生成子と並進生成子の間で零だからである \eqref{eq:killing-degeneracy}。

\paragraph{光円錐上の消滅}
第 2 項 \(\tfrac{1}{2}\eta_{\mu\nu}\lambda^\mu\lambda^\nu\) は、並進パラメータが光円錐 \eqref{eq:light-cone-translation} 上にあるとき正確に\textbf{零}になる：
\begin{equation}
  \eta_{\mu\nu}\lambda^\mu\lambda^\nu = 0
  \quad\Longrightarrow\quad
  J^{(5)} = \tfrac{1}{16}\, B_{\mathrm{Killing}}(\Omega_{\mathrm{torsion}}, \Omega_{\mathrm{torsion}}).
  \label{eq:J5-light-cone}
\end{equation}
この事実は、光的並進が「長さゼロの移動」——ミンコフスキー計量の下で零の二次形式を持つ——として特徴づけられることを意味する。
前章で述べた「光子の移動は並進二次形式が零である」 \eqref{eq:light-cone-translation} という代数的基盤が、\(J^{(5)}\) \eqref{eq:J5-invariant} に組み込まれている。

\paragraph{双対性との整合}
双対写像 \eqref{eq:duality-map} は \(\Omega_{\mathrm{torsion}}\) の双曲成分と巡回成分を入れ替えるが、キリング形式 \eqref{eq:killing-torsion} は \(\mathfrak{so}(3,1) \oplus \mathfrak{so}(3,1)\) 上で適切に定義されているため、\(J^{(5)}\) の第 1 項は双対性の下でも良定義である。
並進係数 \(\lambda^\mu\) はローレンツ 4-ベクトルとして変換されるから、第 2 項もローレンツ共変である。
\(J^{(5)} = 0\) が「完全同期」（ねじれも並進も不整合なし）を意味する条件として、第 IV 部で物理的内容を与えられる。

\section{トーション・ねじれ不整合が現れる代数的場所の準備}

\subsection{10 次元空間上の不整合の座}

PGA 版 DST において、粒子内在の「不整合」はすべて 10 次元バイベクトル空間 \eqref{eq:Omega-decomp} 上の元として記述される。
その\textbf{代数的座}（locus）は次の三箇所に分かれる：

\begin{itemize}[leftmargin=2em]
  \item \textbf{ねじれ 6 次元} \(\Omega_{\mathrm{torsion}}\)：双曲生成子 \(B_a^{+}\) と巡回生成子 \(B_a^{-}\) の線型結合 \eqref{eq:hyperbolic-biv}–\eqref{eq:cyclic-biv}。
  通常時空ローター \(R_{\mathrm{usual}}\) と双対時空ローター \(R_{\mathrm{dual}}\) の相対ずれを担う——第 IV 部で物理的ねじれ不整合として定量化する。
  \item \textbf{並進 4 次元} \(\Omega_{\mathrm{trans}}\)：ヌル生成子 \(N_\mu\) の線型結合 \eqref{eq:null-generators}。
  平面の法線方向に沿った有限シフトを担う——前章のトランスレータ \eqref{eq:exp-truncation} として実装済み。
  \item \textbf{測度 1 次元} \(J^{(5)}\)：上記二セクターを結合するスカラー不変量 \eqref{eq:J5-invariant}。
  ねじれの大きさと並進の大きさを単一の数値で測る。
\end{itemize}

親理論の DST では、重力と慣性は \(\Omega_{\mathrm{torsion}}\) だけで記述されていた。
PGA 版では \(\Omega_{\mathrm{trans}}\) が加わり、光の自由度（零の並進）も同じ 10 次元対象 \eqref{eq:extended-motor} の一部となる——トーション（ねじれ）と並進が別個の概念ではなく、同一代数の異なる成分として等価に扱われる。

\subsection{平面ブラケットによる幾何実現}

前章 \eqref{eq:bracket-def} で導入した平面ベースのブラケット積は、上記の代数的座を幾何的に実現する \eqref{eq:bracket-grades}。

\paragraph{交換子ブラケット：法線＝ヌルバイベクトル}
二つの 3-ブレード \(A, B\) の交換子部分
\begin{equation}
  A \mathbin{\underline{\vee}} B = \langle AB \rangle_2
  \label{eq:bracket-commutator-grade}
\end{equation}
は 2 次部分——バイベクトル——を取り出す。
通常時空の 3-ブレード \(P\) と双対時空要素 \(D\) に対し \eqref{eq:normal-correspondence}、このバイベクトルは
\begin{equation}
  P \mathbin{\underline{\vee}} D = \kappa\,(e_4 \wedge n)
  \label{eq:bracket-normal}
\end{equation}
の形をとり、\(n\) は \(P\) の法線、\(e_4 \wedge n\) はヌルセクター \eqref{eq:null-generators} に属する。
双対時空が「平面の法線束」（normal bundle）として機能する——PGA 版 DST の幾何的核心である。

\paragraph{反交換子ブラケット：並進的ずれ}
反交換子部分
\begin{equation}
  A \mathbin{\overline{\vee}} B = \langle AB \rangle_0 + \langle AB \rangle_4
  \label{eq:bracket-anticommutator-grade}
\end{equation}
は 0 次と 4 次部分を取り出す。
3-ブレード \(P\) に並進パラメータ \(\lambda^\mu\) を含むヌル寄与を加えたとき、
\begin{equation}
  P \mathbin{\overline{\vee}} D
  = \bigl(P \mathbin{\overline{\vee}} D\bigr)_{\mathrm{torsion}}
  + \sum_{\mu=0}^{3} c_\mu\, \lambda^\mu\, \bigl(e_4 \wedge n \wedge e_\mu\bigr),
  \label{eq:bracket-translational}
\end{equation}
と分解される——PGA 版論文 Section 4 の対応する式と一致する。
第 1 項は純粋なねじれ部分（\(\lambda^\mu = 0\) のとき \eqref{eq:Omega-decomp} の \(\Omega_{\mathrm{torsion}}\) に対応）、第 2 項は 4 次ブレード \(e_4 \wedge n \wedge e_\mu\) として並進的不整合 \(\lambda^\mu\) を担う。
したがって \eqref{eq:Omega-decomp} の三分解——ねじれ 6D + 並進 4D——は、平面ベースのブラケット積 \eqref{eq:bracket-def} によって幾何的に実現される。

\subsection{光とねじれの統一的記述}

前章で確立した拡張モーター \(M = \exp(\Omega_{\mathrm{biv}}^{(5)})\) \eqref{eq:extended-motor} は、ねじれローター \(R = \exp(\Omega_{\mathrm{torsion}})\) とトランスレータ \(T = \exp(\Omega_{\mathrm{trans}})\) \eqref{eq:motor-factor} を統合する。
ユニタリ性 \(M\widetilde{M} = 1\) \eqref{eq:motor-unitarity} により、サンドイッチ積 \eqref{eq:motor-sandwich-chapter} は \(G(3,1,1)\) の退化計量を保存する。

この枠組みの下で、光の自由度とねじれ不整合は\textbf{別個の実体}としてではなく、10 次元生成子の異なる成分として等価に扱われる：
\begin{itemize}[leftmargin=2em]
  \item 六つの双曲・巡回生成子 \(\Rightarrow\) 親理論の torsional rotor、ねじれ不整合 \(\Omega_{\mathrm{torsion}}\)
  \item 四つのヌル生成子 \(\Rightarrow\) 並進セクター \(\Omega_{\mathrm{trans}}\)；光円錐 \eqref{eq:light-cone-translation} 上では零
  \item 二次不変量 \(J^{(5)}\) \eqref{eq:J5-invariant} \(\Rightarrow\) 両セクターの結合測度
\end{itemize}
双対写像 \eqref{eq:duality-map} が三分解 \eqref{eq:poincare-iso} を保存する \eqref{eq:null-duality-closure} から、光伝播とねじれ効果は通常／双対セクターの入れ替えの下でも共変に変形する——因果的な零測地線構造が並進不整合の存在下でも維持される。

\subsection{第 IV 部への接続}

本章までで整えた代数的構造は、第 IV 部「二重時空理論の数学的枠組み」の土台となる。

\begin{itemize}[leftmargin=2em]
  \item 通常時空ローター \(R_{\mathrm{usual}}\) と双対時空ローター \(R_{\mathrm{dual}}\) は、\(\Omega_{\mathrm{torsion}}\) の双曲成分と巡回成分として現れる \eqref{eq:duality-hyperbolic-example}。
  \item ねじれ不整合（torsional mismatch）は \(R_{\mathrm{usual}}^\dagger R_{\mathrm{dual}}\) として定義され、\(J^{(5)}\) の第 1 項 \eqref{eq:J5-invariant} がその大きさを測る。
  \item 質量ゼロ粒子（光子）では \(J^{(5)} = 0\) が強制され、特に並進二次形式が零 \eqref{eq:J5-light-cone} となる——光円錐条件 \eqref{eq:light-cone-translation}。
  \item 位置演算子の構成は、通常時空の横方向、双対時空の位相、光的並進の三者構造として第 IV 部で展開される。
\end{itemize}

ここでは物理的定義——\(R_{\mathrm{usual}}\)、\(R_{\mathrm{dual}}\)、ねじれ不整合の定量化——には踏み込まない。
本章の役割は、それらが\textbf{どの代数構造の上}に立つかを明確にすることにある。
第 III 部はこれで完結する——\(G(3,1,1)\) 上の射影構造、モーター、双対性、拡張不変量 \(J^{(5)}\) が整い、DST の数学的枠組みへ進む準備ができた。

\section*{演習問題}

\begin{enumerate}[leftmargin=2em]
  \item \(N_0 = e_4 \wedge e_0 = e_4 j\) \eqref{eq:null-explicit} と \(i = e_0 e_1 e_2 e_3\) \eqref{eq:duality-pseudoscalar} について、\(N_0 i = e_4 (j i)\) を \eqref{eq:e4-i-commute} を用いて計算せよ。
  右辺が \(e_4 \wedge (\text{ベクトル})\) の形——すなわちヌルセクター内——に留まることを確認せよ \eqref{eq:null-duality-closure}。

  \item 双曲生成子 \(iI = e_0 e_1\) \eqref{eq:hyperbolic-biv} に双対写像 \eqref{eq:duality-map} を適用し、\eqref{eq:duality-hyperbolic-example} と同様に \((iI)\, i = K\) となることを確認せよ。
  双曲セクターと巡回セクターが双対の下で入れ替わることを、言葉で述べよ。

  \item キリング形式の退化 \eqref{eq:killing-degeneracy} について、なぜ \(B_{\mathrm{Killing}}(P_\mu, P_\nu) = 0\) であると並進不整合 \(\lambda^\mu\) を検出する二次不変量が得られないのか、2--3 文で説明せよ。

  \item 拡張不変量 \eqref{eq:J5-invariant} の第 2 項について、\(\lambda^\mu = (c, 1, 0, 0)\)（\(c\) は光速）が \eqref{eq:light-cone-translation} を満たすとき \(\tfrac{1}{2}\eta_{\mu\nu}\lambda^\mu\lambda^\nu = 0\) となることを確認せよ。
  このとき \(J^{(5)}\) が \eqref{eq:J5-light-cone} の形に簡約されることを述べよ。

  \item 法線対応 \eqref{eq:normal-correspondence} が双対写像 \eqref{eq:duality-map} の下で \eqref{eq:normal-duality-covariance} の形を保つ理由を、\(N_\mu i = e_4(e_\mu i)\) \eqref{eq:null-duality-closure} を用いてスケッチせよ。

  \item 平面ブラケット \eqref{eq:bracket-def} について、\(\lambda^\mu = 0\) のとき \eqref{eq:bracket-translational} の第 2 項が消え、純粋なねじれ部分だけが残ることを確認せよ。
  これが \eqref{eq:Omega-decomp} の \(\Omega_{\mathrm{trans}} = 0\) の場合に対応することを述べよ。

  \item 双対写像 \eqref{eq:duality-map} と PGA 双対 \eqref{eq:pga-dual} の違いを、作用する対象（セクター vs ブレード次数）と幾何的意味（通常／双対入れ替え vs 点／超平面）の両面から比較せよ。

  \item （発展）純ねじれ（\(\lambda^\mu = 0\)）と純並進（\(\alpha_a = \beta_a = 0\)）の場合に、\(J^{(5)}\) \eqref{eq:J5-invariant} がそれぞれどう簡約されるかを書け。
  第 IV 部で \(J = 0\)（真空）と光子（並進二次形式が零）がどう区別されるかの予習として述べよ。
\end{enumerate}

\newpage

%==================================================
\part{第IV部　二重時空理論の数学的枠組み}

\chapter{通常時空と双対時空}

第 III 部では \(G(3,1,1)\) 上の 10 次元生成子、拡張モーター、双対写像、拡張不変量 \(J^{(5)}\) \eqref{eq:J5-invariant} を整えた。
本章では、DST の物理的核心——各粒子が内在的に持つ\textbf{通常時空}と\textbf{双対時空}——を、第 II 部の \(\mathrm{Cl}(3,1)\) 部分代数と第 III 部の PGA 拡張の言葉で定義する。
背景として一つのミンコフスキー連続体を敷くのではなく、粒子ごとに二つのミンコフスキー構造が共存する——この立場が、以降のねじれ不整合 \eqref{eq:Omega-decomp} と位置演算子の構成の出発点となる。
ねじれの定量化、光円錐条件の強制、位置演算子の厳密な定義は次章以降に委ね、ここでは二つのローター \(R_{\mathrm{usual}}\)、\(R_{\mathrm{dual}}\) とその代数対応を確立する。

\section{\texorpdfstring{二つの時空のローター \(R_{\mathrm{usual}}\) と \(R_{\mathrm{dual}}\)}{二つの時空のローター R usual と R dual}}

\subsection{粒子内在の二重ミンコフスキー構造}

DST の出発点は、時空が宇宙全体に一つだけ存在する連続体である、という従来の発想を改めることにある。
各粒子は、それぞれ\textbf{通常時空}（usual spacetime）と\textbf{双対時空}（dual spacetime）という二つのミンコフスキー構造を内在的に持つ。
両者は \(\mathrm{Cl}(3,1)\) 部分代数 \eqref{eq:cl31-basis} 上の異なる四ベクトル基底として実現され、同一の符号 \((-,+,+,+)\) \eqref{eq:minkowski-signature} の二次形式 \eqref{eq:spacetime-norm} を共有する。

真空や自由落下の理想的状態では、二つの時空は完全に同期している。
質量や運動があると、通常時空側の変換と双対時空側の変換の間にずれが生じる——この相対ずれが\textbf{ねじれ不整合}（torsional mismatch）であり、次章で \(R_{\mathrm{usual}}^\dagger R_{\mathrm{dual}}\) として定量化する。
本章では、ずれが生じる前の「二つの独立なローター」の定義に焦点を当てる。

\subsection{通常時空基底と双対時空基底}

双四元数表記 \eqref{eq:cl31-basis} において、通常時空の四ベクトル基底は
\begin{equation}
  (e_0,\, e_1,\, e_2,\, e_3)
  = (j,\, kI,\, kJ,\, kK)
  \label{eq:usual-basis-dst}
\end{equation}
である。
一般の通常時空ベクトルは
\begin{equation}
  X = ct\, j + x\, kI + y\, kJ + z\, kK = x^\mu e_\mu
  \label{eq:usual-vector-dst}
\end{equation}
と書ける \eqref{eq:spacetime-vector}。

双対時空基底は、時間成分と空間成分の役割分担が入れ替わった第二組の基底として
\begin{equation}
  (e'_0,\, e'_1,\, e'_2,\, e'_3)
  = (k,\, jI,\, jJ,\, jK)
  \label{eq:dual-basis-dst}
\end{equation}
で与えられる。
双対時空ベクトルは
\begin{equation}
  X' = ct'\, k + x'\, jI + y'\, jJ + z'\, jK
  \label{eq:dual-vector-dst}
\end{equation}
と表記する。
両者のミンコフスキーノルムは同一である：
\begin{equation}
  X^2 = X'^2 = -(ct)^2 + x^2 + y^2 + z^2.
  \label{eq:dual-norm-preservation}
\end{equation}
ここで \(ct, x, y, z\) と \(ct', x', y', z'\) は同一粒子の二つの時空座標——必ずしも数値的に等しいわけではなく、二つの内在的自由度として独立に存在する。

\subsection{ローターの定義}

各時空セクターにおけるローレンツ変換は、第 II 部 \eqref{eq:general-lorentz} と同様に偶数ローターで記述される。
通常時空側では、ブースト方向の\textbf{単位双曲バイベクトル} \(\hat{p}\)（\(\hat{p}^2 = +1\)）とラピディティ \(\phi\) を用いて
\begin{equation}
  R_{\mathrm{usual}}
  = \exp\!\Bigl(\tfrac{\phi}{2}\,\hat{p}\Bigr)
  = \cosh\frac{\phi}{2} + \hat{p}\,\sinh\frac{\phi}{2}
  \label{eq:R-usual}
\end{equation}
と定義する。
典型的には \(\hat{p} = p_1\, iI + p_2\, iJ + p_3\, iK\)（\(p_1^2+p_2^2+p_3^2=1\)）であり、双曲生成子 \eqref{eq:hyperbolic-biv} の線型結合である——非コンパクトなブーストを担う。

双対時空側では、回転方向の\textbf{単位巡回バイベクトル} \(\hat{q}\)（\(\hat{q}^2 = -1\)）と位相角 \(\theta\) を用いて
\begin{equation}
  R_{\mathrm{dual}}
  = \exp\!\Bigl(\tfrac{\theta}{2}\,\hat{q}\Bigr)
  = \cos\frac{\theta}{2} + \hat{q}\,\sin\frac{\theta}{2}
  \label{eq:R-dual}
\end{equation}
と定義する。
典型的には \(\hat{q} = q_1 I + q_2 J + q_3 K\)（\(q_1^2+q_2^2+q_3^2=1\)）であり、巡回生成子 \eqref{eq:cyclic-biv} の線型結合である——コンパクトな回転を担う。

\paragraph{物理的分割}
表 \ref{tab:usual-dual-rotors} に、二つのローターの代数的・幾何的対応をまとめる。

\begin{table}[htbp]
  \centering
  \caption{通常時空ローターと双対時空ローターの対応}
  \label{tab:usual-dual-rotors}
  \begin{tabular}{@{}llll@{}}
    \toprule
    & 通常時空 \(R_{\mathrm{usual}}\) & 双対時空 \(R_{\mathrm{dual}}\) \\
    \midrule
    生成子の型 & 双曲（\(\hat{p}^2=+1\)） & 巡回（\(\hat{q}^2=-1\)） \\
    PGA 生成子 & \(iI, iJ, iK\) \eqref{eq:hyperbolic-biv} & \(I, J, K\) \eqref{eq:cyclic-biv} \\
    パラメータ & ラピディティ \(\phi\)（非コンパクト） & 位相 \(\theta\)（コンパクト） \\
    指数展開 & \(\cosh,\sinh\) & \(\cos,\sin\) \\
    物理的役割 & ブースト・外向き運動学 & 回転・内在的位相 \\
    \bottomrule
  \end{tabular}
\end{table}

\subsection{完全ローターと相対ずれ}

二つのローターは、生成元が可換な \(\mathfrak{so}(3,1) \oplus \mathfrak{so}(3,1)\) 構造 \eqref{eq:poincare-iso} の異なる因子に属するため、一般に\textbf{可換}である：
\begin{equation}
  [R_{\mathrm{usual}},\, R_{\mathrm{dual}}] = 0.
  \label{eq:rotor-commute}
\end{equation}
粒子の全体的なローレンツ変換は、両者の積
\begin{equation}
  R_{\mathrm{total}} = R_{\mathrm{usual}}\, R_{\mathrm{dual}}
  \label{eq:R-total}
\end{equation}
で記述される \eqref{eq:spinor-transform}。
スピノル \(\psi\) への作用は \(R_{\mathrm{total}}\) による左乗算 \eqref{eq:spinor-transform}、時空ベクトル \(X\) への作用はサンドイッチ積 \(X' = R_{\mathrm{total}}\, X\, \widetilde{R_{\mathrm{total}}}\) \eqref{eq:lorentz-sandwich} である。

二つのローターの\textbf{相対ずれ}は
\begin{equation}
  \Omega := R_{\mathrm{usual}}^\dagger\, R_{\mathrm{dual}}
  \label{eq:torsion-rotor-preview}
\end{equation}
で測られる——これがねじれ不整合ローターの原型であり、\(\Omega_{\mathrm{torsion}}\) \eqref{eq:Omega-decomp} の双曲成分 \(\alpha_a\) と巡回成分 \(\beta_a\) に対応する \eqref{eq:duality-hyperbolic-example}。
\(J^{(5)}\) \eqref{eq:J5-invariant} の第 1 項は、この相対ずれの大きさを二次形式で測る。

\subsection{PGA 版における第三のセクター}

PGA 版 DST では、並進が \(G(3,1,1)\) のヌル生成子 \(N_\mu = e_4 \wedge e_\mu\) \eqref{eq:null-generators} として内在化される。
並進は通常時空・双対時空の「二番目のミンコフスキー構造」ではなく、射影方向 \(e_4\) から生じる\textbf{第三の自由度}である。
拡張モーター \(M = \exp(\Omega_{\mathrm{biv}}^{(5)})\) \eqref{eq:extended-motor} は、ねじれ部分 \(\Omega_{\mathrm{torsion}}\)（\(R_{\mathrm{usual}}\) と \(R_{\mathrm{dual}}\) のずれ）と並進部分 \(\Omega_{\mathrm{trans}}\)（\(N_\mu\) の線型結合）を統合する。
本章では並進を脇に置き、二つのローターの構造を明確にする——並進の物理的役割は次章で本格的に扱う。

\section{コンパクト化と時間矢印の反転}

\subsection{双対写像による時間反転}

通常時空と双対時空は、第 III 部 \eqref{eq:duality-map} の\textbf{双対写像}
\begin{equation}
  X \;\longmapsto\; X' = X\, i,
  \qquad i = e_0 e_1 e_2 e_3
  \label{eq:duality-map-chapter}
\end{equation}
（\(i\) は 4 次擬スカラー \eqref{eq:cl31-pseudoscalar}）によって結びつく。
基底 \eqref{eq:usual-basis-dst}–\eqref{eq:dual-basis-dst} に対し、成分レベルでは
\begin{align}
  (ct\, j)\, i &= -ct\, k, &
  (x\, kI)\, i &= x\, jI, \notag\\
  (y\, kJ)\, i &= y\, jJ, &
  (z\, kK)\, i &= z\, jK
  \label{eq:duality-basis-action}
\end{align}
となる——時間成分の符号が反転し、空間成分は \((kI,kJ,kK) \mapsto (jI,jJ,jK)\) と写される。
したがって双対時空は、通常時空の\textbf{時間反転鏡像}として読める \eqref{eq:dual-basis-dst}。

ミンコフスキーノルムの保存 \eqref{eq:dual-norm-preservation} は、\((Xi)^2 = X^2 i^2 = -X^2\) より \(X^2 = (Xi)^2\) として確認できる——双対写像は計量を破壊しない。

\paragraph{生成子レベルでの入れ替え}
双対写像は、双曲生成子と巡回生成子を相互に写す \eqref{eq:duality-hyperbolic-example}：
\begin{equation}
  (iI)\, i = K, \qquad (iJ)\, i = J, \qquad (iK)\, i = I.
  \label{eq:duality-generator-swap}
\end{equation}
通常時空のブースト方向 \(\hat{p}\) に対し、
\begin{equation}
  \hat{p}\, i = -\,\hat{q}',
  \qquad
  \hat{q}' := p_1 I + p_2 J + p_3 K,
  \label{eq:boost-to-rotation}
\end{equation}
すなわち双曲バイベクトルは巡回バイベクトルへと変換される \eqref{eq:duality-hyperbolic-example}。
これが「通常時空のブーストが、双対時空では回転として現れる」という主張の代数的内容である。

\subsection{非コンパクト性とコンパクト性}

通常時空側のブーストは\textbf{双曲的}（hyperbolic）である。
ラピディティ \(\phi\) は実数全体を値域とし、\(\phi \to \pm\infty\) で光速に漸近する——\(G(3,1)\) の双曲生成子 \eqref{eq:hyperbolic-biv} は平方 \(+1\) を持ち、指数 \eqref{eq:exp-boost} は \(\cosh,\sinh\) で展開される。
この非コンパクト性が、光子が光速で進むとき進行方向の座標がローレンツ収縮によって「潰れる」現象の代数的原因である \eqref{eq:boost-rotor}。

双対時空側の回転は\textbf{楕円的}（elliptic）——すなわち\textbf{コンパクト}である。
位相角 \(\theta\) は \(4\pi\) 周期 \footnote{スピノル表現では \(2\pi\) 周期だが、DST では双対ローターの位相空間全体が \(SO(3)\) の二重被覆 \(S^3\) として現れる——第 V 部で詳述。} で、\(\theta\) と \(\theta + 4\pi n\)（\(n \in \mathbb{Z}\)）は同一の物理状態を表す。
巡回生成子 \eqref{eq:cyclic-biv} は平方 \(-1\) を持ち、指数 \eqref{eq:exp-rotation} は \(\cos,\sin\) で展開される。
双曲的な「引っ張り」が、双対写像 \eqref{eq:boost-to-rotation} の下で楕円的な「回り」に置き換わる——通常時空で非コンパクトに伸びる自由度が、双対時空では有限の位相として閉じる。

\paragraph{直感的対比}
通常時空のブーストは、一方向に伸び続ける双曲運動に似る。
双対時空の回転は、輪の上を回る運動に似る——速さを上げても「長さ」は変わらず、位相だけが蓄積する。
この対比が、進行方向情報の保存メカニズムの鍵となる（次節・次章）。

\subsection{進行方向位相の保存}

ニュートン‐ウィグナー問題 \eqref{tab:tools-nw} の核心は、質量ゼロ粒子の進行方向座標が通常時空だけでは定義できないことにある。
双対時空は、この失われがちな情報を\textbf{コンパクトな位相}として受け止める役割を担う \eqref{tab:tools-nw}。
通常時空でローレンツブーストにより進行方向が収縮するとき、対応する双対時空の位相角 \(\theta\) が同じ大きさで蓄積する——完全反同期 \eqref{eq:anti-sync} の下では \(\phi = \theta\) となり、情報の損失なく写像される。

PGA 版では、双対時空要素が 3-ブレードの法線として機能する \eqref{eq:normal-correspondence}。
通常時空側の 3-ブレード \(P\) と双対時空セクターの元 \(D\) に対し
\begin{equation}
  P \mathbin{\underline{\vee}} D = \kappa\,(e_4 \wedge n),
  \label{eq:normal-correspondence-chapter}
\end{equation}
ここで \(n\) は \(P\) の法線ベクトル、\(e_4 \wedge n\) はヌルセクター \eqref{eq:null-generators} に属する。
双対時空は、平面幾何の\textbf{法線束}（normal bundle）として読める——PGA 版論文 Section 4 の幾何的核心である \eqref{eq:normal-correspondence}。
法線方向の情報は並進生成子 \(N_\mu\) に符号化され、次章で位置演算子の構成に現れる。

\section{完全反同期の条件（質量ゼロ粒子の場合）}

\subsection{完全反同期の定義}

質量 \(m = 0\) の粒子——とくに光子——では、通常時空と双対時空のローターが\textbf{完全反同期}（perfect anti-synchronisation）という特別な関係を満たす。
すなわち、ラピディティ \(\phi\) と位相 \(\theta\) が等しく、生成子の向きが双対写像 \eqref{eq:duality-map} によって逆向きになる：
\begin{equation}
  \boxed{\;
  \phi = \theta,
  \qquad
  \hat{p}\, i = -\,\hat{q}.
  \;}
  \label{eq:anti-sync}
\end{equation}
第一条件は二つの角度パラメータの一致、第二条件は双曲生成子と巡回生成子が双対写像 \eqref{eq:boost-to-rotation} の下で符号付き対応することを意味する。

\(\hat{p} = p_1 iI + p_2 iJ + p_3 iK\)、\(\hat{q} = q_1 I + q_2 J + q_3 K\) と書けば、\eqref{eq:anti-sync} は
\begin{equation}
  q_a = -p_a \quad (a = 1,2,3), \qquad \phi = \theta
  \label{eq:anti-sync-components}
\end{equation}
と同値である。

\subsection{ねじれ不整合のゼロ固定}

完全反同期 \eqref{eq:anti-sync} の下では、相対ずれローター \eqref{eq:torsion-rotor-preview} が単位元に一致する：
\begin{equation}
  \Omega = R_{\mathrm{usual}}^\dagger R_{\mathrm{dual}} = 1.
  \label{eq:zero-torsion-massless}
\end{equation}
したがってねじれ不整合 \(\Omega_{\mathrm{torsion}} = 0\) \eqref{eq:Omega-decomp} が強制され、拡張不変量 \eqref{eq:J5-invariant} の第 1 項
\begin{equation}
  \tfrac{1}{16}\, B_{\mathrm{Killing}}(\Omega_{\mathrm{torsion}}, \Omega_{\mathrm{torsion}}) = 0
  \label{eq:J5-torsion-zero}
\end{equation}
も零になる。
質量ゼロ粒子では「通常時空と双対時空の角度のズレ」——DST における重力・慣性の代数的原型——が最初から存在しない。

\subsection{物理的帰結}

完全反同期 \eqref{eq:anti-sync} のもとで、光子の状態は次のように理解できる。

\begin{itemize}[leftmargin=2em]
  \item \textbf{通常時空}：進行方向に沿った座標はローレンツブースト \eqref{eq:R-usual} により収縮するが、進行方向に\textbf{垂直}な横方向座標 \((x_\perp, y_\perp)\) は保存される。
  \item \textbf{双対時空}：進行方向に関する情報は位相角 \(\theta\) としてコンパクトに保存される \eqref{eq:R-dual}。\(4\pi\) 周期の位相空間は次章以降 \(S^3\) 構造 \eqref{tab:tools-nw} として量子化される。
  \item \textbf{並進}（予告）：\(J^{(5)} = 0\) と並進パラメータの光円錐制約 \eqref{eq:light-cone-translation} により、実際の移動は純粋な光速並進として実行される \eqref{eq:J5-light-cone}。
\end{itemize}

三者の役割分担を表 \ref{tab:triple-structure} にまとめる。

\begin{table}[htbp]
  \centering
  \caption{PGA 版 DST における三者構造（光子）}
  \label{tab:triple-structure}
  \begin{tabular}{@{}lll@{}}
    \toprule
    セクター & 担う情報 & 代数的実体 \\
    \midrule
    通常時空 & 横方向座標、標準的 SR 現象 & \(R_{\mathrm{usual}}\), \(\{iI,iJ,iK\}\) \\
    双対時空 & 進行方向位相、コンパクト自由度 & \(R_{\mathrm{dual}}\), \(\{I,J,K\}\) \\
    並進 & 光的移動、実際の位置シフト & \(N_\mu\), \(\lambda^\mu\) \eqref{eq:null-generators} \\
    \bottomrule
  \end{tabular}
\end{table}

位置演算子は、この三者を組み合わせて
\[
  \text{位置} = \text{（横方向）} + \text{（双対位相）} + \text{（光的並進）}
\]
として構成される——厳密な定義と共変性の証明は次章「ねじれ不整合と並進の役割」および第 V 部「ニュートン‐ウィグナー問題の解決」に委ねる。

\subsection{質量を持つ粒子との対比}

質量 \(m > 0\) の粒子では、完全反同期 \eqref{eq:anti-sync} は一般に成り立たない。
加速や外部場により \(R_{\mathrm{usual}}\) が変化しても、双対時空のコンパクト剛性 \footnote{双対時空の位相空間のコンパクト性が、質量に比例した「剛性」を生む——親理論 DST の核心的主張。} により \(R_{\mathrm{dual}}\) が即座に追従できず、\(\Omega \neq 1\) となる。
この相対ずれ \(\Omega_{\mathrm{torsion}}\) が、DST において慣性と重力の代数的起源である \eqref{eq:J5-invariant}。
光子は剛性がゼロだから完全反同期が可能であり、質量粒子では剛性がねじれを生む——この対比が、質量の有無による物理の分岐を説明する。

\subsection{次章への接続}

本章で確立した内容を整理する。

\begin{itemize}[leftmargin=2em]
  \item \(R_{\mathrm{usual}}\) \eqref{eq:R-usual} と \(R_{\mathrm{dual}}\) \eqref{eq:R-dual} は、\(\mathfrak{so}(3,1)\) の双曲因子と巡回因子として 10 次元空間 \eqref{eq:poincare-iso} の前半 6 次元に住む。
  \item 双対写像 \eqref{eq:duality-map} は両セクターを入れ替え、時間矢印を反転させる \eqref{eq:duality-basis-action}。
  \item 光子では完全反同期 \eqref{eq:anti-sync} により \(\Omega_{\mathrm{torsion}} = 0\) \eqref{eq:zero-torsion-massless}。
  \item 並進セクター \(N_\mu\) \eqref{eq:null-generators} が第三の自由度として加わり、位置の一意性を完成させる \eqref{tab:triple-structure}。
\end{itemize}

次章では、\(\Omega_{\mathrm{torsion}}\) の厳密な定義、\(J^{(5)} = 0\) から光円錐条件 \eqref{eq:light-cone-translation} が導かれる仕組み、および位置演算子の代数的構成を扱う。

\section*{演習問題}

\begin{enumerate}[leftmargin=2em]
  \item 通常時空ベクトル \(X = ct\, j + x\, kI\) \eqref{eq:usual-vector-dst} に双対写像 \(X \mapsto Xi\) \eqref{eq:duality-map} を適用せよ。
  結果が \eqref{eq:duality-basis-action} の形 \(X' = -ct\, k + x\, jI\) になることを確認し、\(X^2 = X'^2\) \eqref{eq:dual-norm-preservation} も直接計算で検証せよ。

  \item \(\hat{p} = iI\)（\(z\) 方向ブースト）について、\(\hat{p}^2 = +1\) \eqref{eq:hyperbolic-biv} を計算で確認せよ。
  同様に \(\hat{q} = I\) \eqref{eq:cyclic-biv} について \(\hat{q}^2 = -1\) を示せ。

  \item 双曲生成子 \(iI\) \eqref{eq:hyperbolic-biv} に双対写像 \eqref{eq:duality-map} を適用し、\((iI)\, i = K\) \eqref{eq:duality-hyperbolic-example} となることを第 III 部と同様に計算せよ。
  \(\hat{p} = p_1 iI + p_2 iJ + p_3 iK\) に対し \(\hat{p}\, i = -\hat{q}'\) \eqref{eq:boost-to-rotation} が成り立つことを確認せよ。

  \item \(R_{\mathrm{usual}} = \cosh(\phi/2) + iI\,\sinh(\phi/2)\) \eqref{eq:R-usual} と \(R_{\mathrm{dual}} = \cos(\theta/2) + I\,\sin(\theta/2)\) \eqref{eq:R-dual} について、\(\phi = \theta\) かつ \(\hat{p}\, i = -\hat{q}\)（\(\hat{p} = iI\), \(\hat{q} = I\)）のとき \(R_{\mathrm{usual}}^\dagger R_{\mathrm{dual}}\) が \(1\) に近い形になることを、小角度 \(\phi = \theta \ll 1\) の近似で確認せよ \eqref{eq:anti-sync}。

  \item 表 \ref{tab:usual-dual-rotors} を参照し、通常時空のブーストが「非コンパクト」、双対時空の回転が「コンパクト」である理由を、生成子の平方 \((+1)\) vs \((-1)\) と指数展開 \((\cosh,\sinh)\) vs \((\cos,\sin)\) の観点から 3--4 文で説明せよ。

  \item 法線対応 \eqref{eq:normal-correspondence} について、双対時空要素 \(D\) が 3-ブレード \(P\) の法線として機能することの幾何的意味を、平面とその法線方向の直感を用いて述べよ。
  交換子ブラケット \(P \mathbin{\underline{\vee}} D\) がヌルバイベクトル \(e_4 \wedge n\) \eqref{eq:normal-correspondence-chapter} になることと、並進生成子 \(N_\mu\) \eqref{eq:null-generators} との関係を説明せよ。

  \item 完全反同期 \eqref{eq:anti-sync} が成り立つとき、\(J^{(5)}\) \eqref{eq:J5-invariant} の第 1 項 \eqref{eq:J5-torsion-zero} が零になる理由を、\(\Omega_{\mathrm{torsion}} = 0\) \eqref{eq:zero-torsion-massless} から説明せよ。
  第 2 項 \(\tfrac{1}{2}\eta_{\mu\nu}\lambda^\mu\lambda^\nu\) が光円錐 \eqref{eq:light-cone-translation} 上で零になること \eqref{eq:J5-light-cone} との関係を、次章の予習として 2--3 文で述べよ。

  \item （発展）\(R_{\mathrm{total}} = R_{\mathrm{usual}} R_{\mathrm{dual}}\) \eqref{eq:R-total} について、\([R_{\mathrm{usual}}, R_{\mathrm{dual}}] = 0\) \eqref{eq:rotor-commute} が成り立つ理由を、生成子が \(\mathfrak{so}(3,1) \oplus \mathfrak{so}(3,1)\) の異なる因子に属すること \eqref{eq:poincare-iso} から説明せよ。
  質量 \(m > 0\) の粒子で \(\Omega \neq 1\) \eqref{eq:torsion-rotor-preview} となる物理的意味を、双対時空の「剛性」と関連づけて述べよ。
\end{enumerate}

\newpage

\chapter{ねじれ不整合と並進の役割}

\section{ねじれ不整合（torsional mismatch）の定義}
% 通常時空と双対時空のずれの定量化。

\section{光子における光円錐条件の強制}
% ねじれ不整合がゼロのときの帰結。

\section{位置演算子の構成}
% 通常時空の横方向 ＋ 双対時空の位相 ＋ 光的並進。

\section*{演習問題}

\chapter{スピノル・量子状態との統合}

\section{双対時空ローターとディラックスピノルの関係}
% スピノル表現との自然な接続。

\section{\texorpdfstring{ヒルベルト空間の構成（特に \(S^3\) 構造）}{ヒルベルト空間の構成（特に S3 構造）}}
% コンパクト位相空間による量子状態の表現。

\section{モーターによる位置とスピンの統一変換}
% サンドウィッチ積による同時変換。

\section*{演習問題}

\newpage

%==================================================
\part{第V部　応用と展望}

\chapter{ニュートン‐ウィグナー問題の解決}

\section{質量ゼロ粒子の局所化がなぜ可能になるか}
% 三者構造による情報の保存と復元。

\section{観測者独立性の回復}
% 位置演算子の共変性。

\section*{演習問題}

\chapter{重力・慣性への拡張の萌芽}

\section{ねじれ不整合としての重力・慣性}
% 重力の幾何代数的解釈への入口。

\section{量子重力への道筋（概観）}
% 今後の研究課題の展望。

\section*{演習問題}

\newpage

%==================================================
\appendix

\chapter{記法一覧と符号規約}
% 本書で用いる記号・符号の統一表。

\chapter{主要な同型関係}
% 四元数・双四元数・ローターなどの関係をまとめる。

\chapter{\texorpdfstring{\(G(3,1,1)\) の基底と積表}{G(3,1,1) の基底と積表}}
% 計算用チートシート。

\chapter{参考文献と発展的トピック}
% デイビッド・ヘステネス、チャールズ・ガン、エリック・レンジェルなどの主要文献。
% その他、関連する幾何代数・相対論・量子重力の文献。

\end{document}
